\section{Shared results VIII: further shared results}\label{sec:pf:sh-misc}

\providecommand{\shmiscA}{\mathcal{A}}
\providecommand{\shmiscC}{\mathcal{C}}
\providecommand{\shmiscFm}{\mathcal{F}}
\providecommand{\shmiscLk}{\mathrm{Lk}}
\providecommand{\shmiscKP}{\mathrm{KP}}
\providecommand{\shmiscBd}{\mathrm{Bd}}
\providecommand{\shmiscQO}{\mathrm{QO}}
\providecommand{\shmiscand}{\mathbin{\wedge}}
\providecommand{\shmiscEv}{\mathrm{Ev}}
\providecommand{\shmiscLg}{\mathrm{Lg}}
\providecommand{\shmiscBu}{\mathrm{Bu}}
\providecommand{\shmiscLin}{\mathrm{Lin}}
\providecommand{\shmiscLB}{\mathrm{LB}}
\providecommand{\shmiscE}{\mathcal{E}}
\providecommand{\shmiscTrue}{\mathrm{True}}
\providecommand{\shmiscIrr}{\operatorname{Irr}}
\providecommand{\shmiscKone}{\mathrm{K}_1}
\providecommand{\shmiscGo}{\ensuremath{(\mathrm{G}^\circ)}}
\providecommand{\shmiscS}{\mathcal S}
\providecommand{\shmiscR}{\mathcal R}
\providecommand{\shmiscX}{\mathsf X}
\providecommand{\shmiscXO}{\mathsf X_{\OO}}
\providecommand{\shmiscla}{\lambda}

This section proves twenty-eight results of the formal development which are used by several sections of Part~II and which
are not proved there. They are of very different weight, and we group them by topic.
\begin{enumerate}[label=(\arabic*)]
\item \emph{Rows and eventual statements} (\cref{sec:pf:sh-misc:rows}): seven formal properties of the eventual statements, rows and
linked rows of \cref{def:families:rows} (trivial rows, monotonicity, transport, conjunction, the trivial large-base condition, and the
duties of designed classes), used in \cref{sec:pf:rows,sec:pf:stageU,sec:pf:thmS,sec:pf:open0,sec:pf:reckernel,sec:pf:baselat1}.
\item \emph{Standing hypotheses and Theorem~B} (\cref{sec:pf:sh-misc:standing}): the hypotheses of the $K_2$-theorems and the standard
conditions follow from the standing hypotheses; the global Theorem~B for the mass data is assembled from four uniform statements; and the
conditions on complex data at the base level follow from the primitive links of a point. These are used in
\cref{sec:pf:open0,sec:pf:thmS,sec:pf:stageU,sec:pf:fit,sec:pf:lift2,sec:pf:onto2,sec:pf:onto3}.
\item \emph{Hosts} (\cref{sec:pf:sh-misc:hosts}): elementary facts on the level, the characters, the bad primes and the case division
of the hosts of the window bound, and on coefficient rings (\cref{sec:pf:vwin,sec:pf:eis,sec:pf:modularity}).
\item \emph{Host frames over $\Q(i)$} (\cref{sec:pf:sh-misc:frames}): the Galois group of $\Q(i)$ and condition $(\mathrm{F1})$ for host
frames (\cref{sec:pf:eis,sec:pf:modularity}).
\item \emph{Smooth representations of $\GL_n$} (\cref{sec:pf:sh-misc:llc}): depth zero and tame twists, goodness and supercuspidal
supports (\cref{sec:pf:llc,sec:pf:galred}).
\item \emph{Class field theory and $L$-functions} (\cref{sec:pf:sh-misc:cft}): ray class fields from global Artin reciprocity, partial
Hecke $L$-functions, and the integrality of singular moduli
(\cref{sec:pf:baselat2,sec:pf:onto3,sec:pf:bundles,sec:pf:liftinj,sec:pf:reckernel,sec:pf:rows}).
\item \emph{Intersection numbers} (\cref{sec:pf:sh-misc:intnum}): the intersection numbers defined by multigraded Hilbert functions form a
field intersection family, and such families are unique on projective varieties (\cref{sec:pf:deepstr,sec:pf:towergeo}).
\item \emph{The second Lie derivative formula} (\cref{sec:pf:sh-misc:lie2}): the infinitesimal action of the second unitary group of a
Weil representation on the Fock model, with the scalar shift $\frac{m_w}2\tr Y_w$ determined by a doubling argument
(\cref{sec:pf:thmK,sec:pf:tlift}).
\item \emph{Howe's unramified generation} (\cref{sec:pf:sh-misc:howe}): in a lattice model of an unramified local dual pair of unitary
groups, the vectors fixed by the integral points of one group are generated by the vacuum under the other
(\cref{sec:pf:bundles,sec:pf:liftinj,sec:pf:open0,sec:pf:stageU}).
\item \emph{Norm-one primes in short intervals} (\cref{sec:pf:sh-misc:shortnorm}): a lower bound, uniform in the discriminant and the
modulus, for the primes $p=\pi\bar\pi$ with $\pi\equiv1\bmod M$ in $\Q(\sqrt{-d})$ in intervals $(X,X(1+h)]$
(\cref{sec:pf:stageU,sec:pf:top-supply}).
\end{enumerate}
Results (1)--(7) are proved completely, from results proved elsewhere in Part~II (cited by their labels) and standard facts. For (8)--(10),
which are the three substantial results of the section, we give the complete architecture of the formal proofs and the full arguments at the
top levels: the doubling argument which determines the scalar in (8); the lattice model, the induction on the index of the classes and the
reduction to two statements on integral Gram matrices in (9); the reduction of (10) to an axiomatic short-interval theorem for
``$L$-families'', and the deduction of that theorem from four analytic inputs (zero-free region, uniqueness and effective separation of an
exceptional zero, and an explicit formula). The deeper ingredients of (8)--(10) are stated precisely and cited by their formal names; they
are proved in the formal development \cite{LeanSource2026}, and their natural-language proofs are not given in this section.

% ===========================================================================
\subsection{Rows and eventual statements}\label{sec:pf:sh-misc:rows}

This subsection proves seven formal properties of the eventual statements and of the rows of
\cref{def:families:rows}, which are used in the sections on rows, on the non-emptiness of the hypothesis package, on
Theorem~S and on the open statements (\cref{sec:pf:rows,sec:pf:stageU,sec:pf:thmS,sec:pf:open0,sec:pf:reckernel,sec:pf:baselat1}).
We use the definitions of \cref{pf:stageU:ev,pf:stageU:rows,pf:stageU:linked} and their notation: the eventual
statements $\shmiscEv_K$, $\shmiscEv^{\mathrm g}_K$ and $\shmiscEv_{\kappa,K}$, the statement $\shmiscLg_K$ ``for large
twisting data'', bundles $\shmiscBu_\kappa(Q,\shmiscLk)$ built from a slot $Q$ along a linkage $\shmiscLk$ at a
parameter tuple $\kappa$, the statement $\shmiscE(Q,\shmiscLk;E)$ that a bundle property $E$ holds eventually at linked
bundles, $E$-selections, the rows $\sfname{Row}(Q,\shmiscLk,H)$, the large-base condition
$\shmiscLB_{\shmiscKP}(Q,R)$ and the linked rows $\shmiscLin_{\shmiscKP}(Q,Q';\shmiscLk,H)$; $Q\shmiscand R$ is the
pointwise conjunction of two slots, and $\shmiscTrue$ denotes the slot (or the bundle property) which holds everywhere.
For a property $\shmiscQO$ of data with objects, $\shmiscC[\shmiscQO]$ is the class of families of
\cref{pf:stageU:rows} (the paragraph before it): families satisfying the page conditions, all of whose members
$(D_{x,i},O_{x,i})$, $x>x_1$, satisfy $\shmiscQO$, the condition $\mathrm{QD}$ and the point conditions.

Two elementary features of the eventual statements are used repeatedly. (i) \emph{Monotonicity}: if $p(K)$ implies
$q(K)$ for every context $K$, then $\shmiscEv_K\,p$ implies $\shmiscEv_K\,q$ (keep all the witnesses); the same holds for
$\shmiscEv^{\mathrm g}_K$, for $\shmiscEv_{\kappa,K}$ (where it suffices that $p(\kappa,K)\Rightarrow q(\kappa,K)$ whenever
$\kappa\triangleright K$), and for $\shmiscLg_K$. (ii) \emph{Conjunction}: two eventual statements of the same kind imply
the eventual statement of the conjunction. For $\shmiscEv_K$ one takes the minimum of the two $\eps_0$, the supremum of the two
transported constants $K_0,K_0'$ (a constant package which is $(K_0\vee K_0')$-large is $K_0$-large and
$K_0'$-large), and the maxima of the thresholds $c_1$, $c_2$, $x_0$; for $\shmiscLg_K$ one takes the unions of the
finite sets $S_1,P_1,\mathcal E(\ell),S_2$ (\lref{GlobalTheory.Mem.Top.large₂.inter}, which states this for
$\shmiscLg_K$); see also \lref{GlobalTheory.Mem.Top.EvG.inter} for $\shmiscEv^{\mathrm g}_K$.

\begin{lemma}[Trivial eventual statements]\label{pf:sh-misc:evtriv}
\begin{enumerate}[label=(\alph*)]
\item For every property $p$ of contexts, $\shmiscEv_K\,p(K)$ implies $\shmiscEv^{\mathrm g}_K\,p(K)$.
\item For every base predicate $\shmiscKP$ and every slot $Q$, the large-base condition $\shmiscLB_{\shmiscKP}(Q,\shmiscTrue)$
holds.
\end{enumerate}
\formallabel{GlobalTheory.Mem.Top.evG_of_ev}\formallabel{GlobalTheory.Fc.largeBaseProperty_true}
\end{lemma}

\begin{proof}
(a) Let $\eps_0$ witness $\shmiscEv_K\,p$. Given $\eps,C_8$ and the three further constants $C_5,c_6,C_b$ of
$\shmiscEv^{\mathrm g}_K$, take the transported constant $K_0$ provided by $\shmiscEv_K\,p$ for $(\eps,C_8)$, and for every
admissible constant package and all $c_D,c_h$ the thresholds $c_1,c_2,x_0$ provided by $\shmiscEv_K\,p$. These witness
$\shmiscEv^{\mathrm g}_K\,p$: the additional hypothesis of $\shmiscEv^{\mathrm g}_K$, that $K$ is $(C_5,c_6,C_b)$-good, is
not used.

(b) By definition, $\shmiscLB_{\shmiscKP}(Q,\shmiscTrue)$ is $\shmiscEv_K\,\shmiscLg_K\,p_K$, where $p_K(B,f,T)$ states that
if $\shmiscKP(B,f)$, then every core over stages over $B$ which satisfies $Q$ satisfies $\shmiscTrue$. Thus $p_K$ holds
for all $(B,f,T)$, and $\shmiscLg_K\,p_K$ holds for every context $K$, with all four exceptional sets $S_1,P_1,\mathcal E(\ell),S_2$
empty (the requirement that their elements be algebraic numbers is then vacuous). Hence $\shmiscEv_K\,\shmiscLg_K\,p_K$ holds
with $\eps_0=1$, an arbitrary transported constant $K_0$ (the formal proof takes the zero constant) and the thresholds
$c_1=c_2=x_0=0$.
\end{proof}

\begin{lemma}[Rows from selections; monotonicity of rows]\label{pf:sh-misc:rowsel}
Let $Q$ be a slot and $\shmiscLk$ a linkage for a bundle type $\shmiscBd$, and let $H,H'$ be properties of classes with raw
selections.
\begin{enumerate}[label=(\alph*)]
\item If $H(\shmiscA,\sigma)$ holds for every class $\shmiscA$ and every $\shmiscTrue$-selection $\sigma$ from $(Q,\shmiscLk)$
for $\shmiscA$, then $\sfname{Row}(Q,\shmiscLk,H)$.
\item If $\sfname{Row}(Q,\shmiscLk,H)$, and $H(\shmiscA,\sigma)$ implies $H'(\shmiscA,\sigma)$ for every class $\shmiscA$ and
every $\shmiscTrue$-selection $\sigma$ from $(Q,\shmiscLk)$ for $\shmiscA$, then $\sfname{Row}(Q,\shmiscLk,H')$.
\end{enumerate}
\formallabel{GlobalTheory.Fc.Row.of_sel}\formallabel{GlobalTheory.Fc.Row.imp}
\end{lemma}

\begin{proof}
(a) Take for $E$ the bundle property $\shmiscTrue$. The statement $\shmiscE(Q,\shmiscLk;\shmiscTrue)$ is of the form
$\shmiscEv_{\kappa,K}$ of a statement which holds for all $\kappa$ and $K$; it holds with $\eps_0=1$, an arbitrary
transported constant, $c_1=c_2=0$ and, for every positive $\kappa$, $x_0=0$. An $E$-selection for $E=\shmiscTrue$ is a
$\shmiscTrue$-selection, so $H$ holds for it by hypothesis.

(b) Let $E$ witness $\sfname{Row}(Q,\shmiscLk,H)$. Every $E$-selection $\sigma$ from $(Q,\shmiscLk)$ for $\shmiscA$ is a
$\shmiscTrue$-selection (forget the condition $E(\kappa,\cdot)$ on the selected bundles); hence $H(\shmiscA,\sigma)$ holds
by the choice of $E$, and $H'(\shmiscA,\sigma)$ by hypothesis. So $E$ witnesses $\sfname{Row}(Q,\shmiscLk,H')$.
\end{proof}

\begin{lemma}[Transport of rows]\label{pf:sh-misc:rowmap}
Let $\shmiscBd,\shmiscBd'$ be bundle types and $\pi=(\pi_{D,O,c})$ a family of maps $\pi_{D,O,c}\colon\shmiscBd(D,O,c)\to
\shmiscBd'(D,O,c)$. Let $Q,Q'$ be slots with $Q\subseteq Q'$ (every core satisfying $Q$ satisfies $Q'$), $\shmiscLk$ a
linkage for $\shmiscBd$ and $\shmiscLk'$ one for $\shmiscBd'$ such that $\shmiscLk(K,\beta,c,b)$ implies
$\shmiscLk'(K,\beta,c,\pi(b))$ for every bundle point $\beta$ over $K$ in $Q$. For a raw selection $\sigma$ with values in
$\shmiscBd$ let $\pi\circ\sigma$ be the raw selection $(\shmiscFm,x,i)\mapsto\pi(\sigma(\shmiscFm,x,i))$. Then for every
property $H$ of classes with raw selections in $\shmiscBd'$,
\[
\sfname{Row}(Q',\shmiscLk',H)\ \Longrightarrow\ \sfname{Row}\bigl(Q,\shmiscLk,(\shmiscA,\sigma)\mapsto H(\shmiscA,\pi\circ\sigma)\bigr).
\]
\formallabel{GlobalTheory.Fc.Row.map}
\end{lemma}

\begin{proof}
Let $E$ witness $\sfname{Row}(Q',\shmiscLk',H)$ and put $E^\pi(\kappa,D,O,c,b)=E(\kappa,D,O,c,\pi(b))$.

\emph{$E^\pi$ holds eventually at linked bundles.} If $\beta$ is a bundle point over $K$ in $Q$ and $\shmiscLk(K,\beta,c,b)$, then
$\beta$ lies in $Q'$ and $\shmiscLk'(K,\beta,c,\pi(b))$; so the statement inside $\shmiscE(Q',\shmiscLk';E)$ implies, for
every $(\kappa,K)$, the statement inside $\shmiscE(Q,\shmiscLk;E^\pi)$, and $\shmiscE(Q,\shmiscLk;E^\pi)$ follows by monotonicity.

\emph{Built bundles.} If $b\in\shmiscBu_\kappa(Q,\shmiscLk)$, witnessed by $K$, a bundle point $\beta$ in $Q$ and $b'$ with
$\shmiscLk(K,\beta,c,b')$ and $(D,O,b)=(\beta.D,\beta.O,b')$, then the same $K,\beta$ and $\pi(b')$ show
$\pi(b)\in\shmiscBu_\kappa(Q',\shmiscLk')$: indeed $\beta$ lies in $Q'$, $\shmiscLk'(K,\beta,c,\pi(b'))$, and applying $\pi$ to the
equality of triples gives $(D,O,\pi(b))=(\beta.D,\beta.O,\pi(b'))$.

\emph{Selections.} Let $\sigma$ be an $E^\pi$-selection from $(Q,\shmiscLk)$ for a class $\shmiscA$. For $\shmiscFm\in\shmiscA$
take the parameter tuple $\kappa\sim\shmiscFm$ given by $\sigma$; then $\shmiscFm\in\shmiscC[\shmiscTrue]$, every
$\pi(\sigma(\shmiscFm,x,i))$ is built from $Q'$ along $\shmiscLk'$ at $\kappa$ by the previous paragraph, and it satisfies
$E(\kappa,\cdot)$ because $\sigma(\shmiscFm,x,i)$ satisfies $E^\pi(\kappa,\cdot)$. Hence $\pi\circ\sigma$ is an $E$-selection from
$(Q',\shmiscLk')$ for $\shmiscA$, and $H(\shmiscA,\pi\circ\sigma)$ holds. Thus $E^\pi$ witnesses the transported row.
\end{proof}

\begin{lemma}[Conjunction of linked rows]\label{pf:sh-misc:rowLand}
Let $\shmiscKP$ be a base predicate, $Q,Q'$ slots, $\shmiscLk$ a linkage and $H,H'$ properties of classes with raw
selections. If $\shmiscLin_{\shmiscKP}(Q,Q';\shmiscLk,H)$ and $\shmiscLin_{\shmiscKP}(Q,Q';\shmiscLk,H')$, then
$\shmiscLin_{\shmiscKP}(Q,Q';\shmiscLk,H\wedge H')$.
\formallabel{GlobalTheory.Fc.RowL.and}
\end{lemma}

\begin{proof}
Let $R$ and $R'$ be slots with $\shmiscLB_{\shmiscKP}(Q,R)$, $\sfname{Row}(Q'\shmiscand R,\shmiscLk,H)$, $\shmiscLB_{\shmiscKP}(Q,R')$
and $\sfname{Row}(Q'\shmiscand R',\shmiscLk,H')$. The two large-base conditions are statements $\shmiscEv_K\,\shmiscLg_K$; by
the conjunction property of $\shmiscEv_K$ and of $\shmiscLg_K$ and monotonicity, they give $\shmiscLB_{\shmiscKP}(Q,R\shmiscand R')$:
a core satisfying $Q$ satisfies $R$ and $R'$. Since $Q'\shmiscand(R\shmiscand R')$ is contained in $Q'\shmiscand R$ and in
$Q'\shmiscand R'$, the two rows hold for the slot $Q'\shmiscand(R\shmiscand R')$ (a row for a slot is a row for every smaller
slot, \lref{GlobalTheory.Fc.Row.strengthen}), and two rows for the same slot and linkage give the row for the conjunction
(\lref{GlobalTheory.Fc.Row.and}). Hence $R\shmiscand R'$ witnesses $\shmiscLin_{\shmiscKP}(Q,Q';\shmiscLk,H\wedge H')$.
\end{proof}

\begin{lemma}[Duties of designed classes]\label{pf:sh-misc:duties}
Let $\shmiscQO$ be a property of data with objects and $\shmiscFm\in\shmiscC[\shmiscQO]$. Then $\shmiscFm$ fulfils the seven
duties: the two page conditions, the global twisting fact at all members, distinct tame primes for $x\ge x_1+1$, the
walls at all members, and, for $x>x_1$, every tuple of every member $D_{x,i}$ is used and $D_{x,i}$ is twisted at all places with
good twists (with respect to the extra data $\mathrm{ex}(x,i)$). Moreover the lattices $\widetilde M(\mu)$ of the members
$D_{x,i}$, $x\ge x_1+1$, are stable under multiplication by $\OO_{\mathit{LL}}$.
\formallabel{GlobalTheory.EX.DQ.familyClass_sevenDuties_and_stableLattices}
\end{lemma}

\begin{proof}
By definition $\shmiscC[\shmiscQO]$ consists of the families satisfying the two page conditions all of whose members
$(D_{x,i},O_{x,i})$, $x>x_1$, satisfy the point conditions, among them the duties of the point (the global twisting fact
and the tameness condition), the wall condition, $\shmiscQO$ and $\mathrm{QD}$. The five duties of the point conditions
(page conditions, global twisting, tameness, walls) are read off from these conditions; for the tameness condition the
formal statement requires distinct tame primes for $x\ge x_1+1$, which holds since such $x$ exceed $x_1$. The condition
$\mathrm{QD}$ at $(D_{x,i},O_{x,i})$ consists of three parts: every tuple of $D_{x,i}$ is used; $D_{x,i}$ is twisted at all
places with good twists for the twisting data of $O_{x,i}$; the lattices $\widetilde M(\mu)$ are $\OO_{\mathit{LL}}$-stable.
The first two give the remaining two duties with threshold $x_0=x_1$, and the third gives the stability with threshold
$x_0=x_1+1$ (the statement quantifies over all $x\ge x_0$, and $x\ge x_1+1$ implies $x>x_1$).
\end{proof}

% ===========================================================================
\subsection{Standing hypotheses, the global Theorem B and base-level conditions}\label{sec:pf:sh-misc:standing}

This subsection proves four results on the hypotheses of the counting theorems of Theorem~S and of the open
statements: two consequences of the standing hypotheses (used in \cref{sec:pf:open0,sec:pf:thmS}), the global
Theorem~B for the mass data (used in \cref{sec:pf:open0,sec:pf:stageU}), and the derivation of the conditions on the
complex data at the base level from the primitive links of a point (used in
\cref{sec:pf:fit,sec:pf:lift2,sec:pf:onto2,sec:pf:onto3,sec:pf:open0}). The last two are assemblies of results proved
elsewhere.

We use the standing hypotheses $\mathrm{SH}$, the standing numerical data $\mathrm{SN}(D,C_0)$ and
$\mathrm{H}_0(C_0)=\mathrm{SH}\wedge\mathrm{SN}(D,C_0)\wedge(1\le C_0)$ of \cref{pf:open0:def:standing}, for a datum $D$ with
size data $L_2$ (with $L=L_2.L$, $b=e^L$) and constants $(C',C_l,C_g,c_b,C_u,c_T,C_T)$. We write $k,m,f,t$ for the
ranks of $D$ ($k,m$ even, $m\ge2$, $k\ge m+3\ge5$), $r=k+m-1$, $r_S=m/2-1$, $g_0,g_1$ for the two type primes of the host,
$E\supset F$ for the CM field of the host and its totally real subfield, $\mathrm{rd}_E$ for its root discriminant, $T_u$ for the
level prime of a word $u$ (a prime, a field of the datum), and, for an embedding type $\tau$, $\lambda_\tau\in\R^r$ for the list
$(v_1,\dots,v_{k-1},-b_1,\dots,-b_m)$ formed from the Harish-Chandra parameter of $D$ at $\tau$. We use the following conjuncts
of these hypotheses (numbered as in the formal definitions):
\begin{itemize}
\item from $\mathrm{SH}$: (a2) $C'm^5fL\le k$; (a3) $m^4\le f$; (a5) $\log f\le\sqrt{r_S}$; (a11) $\log\mathrm{rd}_E\le L/C_l$;
(a12)--(a15) $(kb)^8\le g_i\le b^{C_g}$ for $i=0,1$; (a22) $C_lkm^4+rm^{10}<\lambda_\tau(i)-\lambda_\tau(i+1)$ for all $\tau$ and
$1\le i<r$; (a23) $b^{c_Tm}\le T_u\le b^{C_Tm}$ for all words $u$; (a24) $C_l\le m$; (a25) $0<C_l$;
\item from $\mathrm{SN}(D,C_0)$: (d9) $T_u$ splits completely in $E$ and $T_u\equiv1\bmod(g_0g_1)^2$ for every word $u$;
(d11), (d12) $g_0$ and $g_1$ split completely in $E$; (d14) $T_u\ne g_0,g_1$ for every word $u$.
\end{itemize}
Moreover $1<L$, the type primes $g_0\ne g_1$ are primes, and $E$ satisfies the condition $\sfname{IsCMUnram}$ ($E$ is totally
complex, the involution of $E$ is complex conjugation under every embedding, and for every maximal ideal $P$ of $\OO_E$ there
is $\theta\in\OO_E$ with $\theta-\bar\theta\notin P$, i.e.\ $E/F$ is unramified at all finite places); these are part of the
datum.

\begin{definition}[$K_2$-hypotheses; standard conditions]\label{pf:sh-misc:k2def}
\begin{enumerate}[label=(\alph*)]
\item The \emph{parameter conditions of the $K_2$-theorems} $\mathrm{P}_{K_2}(D,L_2;C_l,C_g,C_T)$ are: (P1)
$C_lm^6L\le k\log k$; (P2) $(kb)^8\le g_i\le b^{C_g}$ for $i=0,1$; (P3) $C_lkm^3<|\lambda_\tau(a)-\lambda_\tau(c)|$ for every $\tau$
and all $a\ne c$; (P4) $C_l\log\mathrm{rd}_E\le L$; (P5) $\log f\le r_S^2$; (P6) $T_u\le b^{C_Tm}$ and $\gcd(T_u,g_0g_1)=1$ for every
word $u$; (P7) $\sfname{IsCMUnram}(E)$. The \emph{$K_2$-hypotheses} are $\mathrm{H}_0(C_0)\wedge\mathrm{P}_{K_2}(D,L_2;C_l,C_g,C_T)$.
\item For words $u,u'$ of $D$, the \emph{standard conditions} $\mathrm{St}(D,u,u')$ are: $r_S\ge1$; $T_u,T_{u'}$ are primes
$\ge3$ which split completely in $E$; $g_0\ne g_1$ are primes which split completely in $E$; with $g_S=g_0g_1$,
$\gcd(g_S,T_uT_{u'})=1$ and $T_u\equiv T_{u'}\bmod g_S$.
\end{enumerate}
\formal{Formalism.MainNotions.NonGenericAndKTwoTheorems.KTwoHypotheses}
\formal{Formalism.MainNotions.NonGenericAndKTwoTheorems.KTwoParameterConditions}
\formal{Formalism.MainNotions.S.Cnt.Stand}
\end{definition}

These are the hypotheses of the $K_2$-theorems used in \cref{pf:open0:price} and the standard conditions of
\cref{sec:pf:thmS}.

\begin{lemma}[The $K_2$-hypotheses]\label{pf:sh-misc:hypof}
Let $D$ be a datum with size data $L_2$ and let $C',C_l,C_g,c_b,C_u,c_T,C_T,C_0$ be reals with $C'\ge1$ and $C_0\ge1$. If
$\mathrm{SN}(D,C_0)$ and the standing hypotheses $\mathrm{SH}$ hold for these constants, then the $K_2$-hypotheses hold.
\formallabel{Formalism.MainNotions.NonGenericAndKTwoTheorems.Hyp_of}
\end{lemma}

\begin{proof}
$\mathrm{H}_0(C_0)$ is the conjunction of the hypotheses. We verify (P1)--(P7).
(P1) Since $C_l\le m$ by (a24), $L\ge0$, $m\ge2$, (a3), $C'\ge1$ and (a2),
\[
C_lm^6L\le m^7L\le m^9L=m^5\cdot m^4\cdot L\le m^5fL\le C'm^5fL\le k\le k\log k,
\]
the last step because $k\ge5$ gives $\log k\ge1$.
(P2) is (a12)--(a15).
(P3) Fix $\tau$ and put $G=C_lkm^3$; $G\ge0$ by (a25). Since $m\ge1$, $G\le C_lkm^4\le C_lkm^4+rm^{10}$, so (a22) gives
$G<\lambda_\tau(i)-\lambda_\tau(i+1)$ for all $1\le i<r$. Summing consecutive differences, $G<\lambda_\tau(i)-\lambda_\tau(j)$ for all
$i<j$ (by induction on $j-i$, using $G\ge0$), hence $G<|\lambda_\tau(a)-\lambda_\tau(c)|$ for all $a\ne c$.
(P4) follows from (a11) after multiplication by $C_l>0$.
(P5) By (a5), $\log f\le\sqrt{r_S}\le r_S^2$ (for an integer $r_S\ge0$).
(P6) The first inequality is part of (a23); since $T_u$, $g_0$, $g_1$ are primes and $T_u\ne g_0,g_1$ by (d14), $T_u$ is prime to
$g_0g_1$.
(P7) is a property of the datum.
\end{proof}

\begin{lemma}[The standard conditions]\label{pf:sh-misc:standof}
Let $D$, $L_2$ and the constants be as in \cref{pf:sh-misc:hypof} (with arbitrary $C'$ and $C_0$), assume
$\mathrm{SN}(D,C_0)$ and $\mathrm{SH}$, and assume $C_l>2$. Then $\mathrm{St}(D,u,u')$ holds for all words $u,u'$ of $D$.
\formallabel{Formalism.MainNotions.Stand_of}
\end{lemma}

\begin{proof}
\emph{$r_S\ge1$.} By (a24), $m\ge C_l>2$, so $m\ge3$; as $m$ is even, $m\ge4$ and $r_S=m/2-1\ge1$.
\emph{Level primes.} $T_u$ is a prime. If $T_u=2$, then (d9) gives $2\equiv1\bmod(g_0g_1)^2$; but $g_0,g_1\ge2$, so
$(g_0g_1)^2\ge16$ and $2\not\equiv1$; hence $T_u\ge3$. $T_u$ splits completely in $E$ by (d9). The same holds for $u'$.
\emph{Type primes.} $g_0\ne g_1$ are primes (data of the host) and split completely in $E$ by (d11), (d12).
\emph{Coprimality and congruence.} As in (P6), each $T_w$ ($w=u,u'$) is a prime different from $g_0,g_1$ by (d14), hence prime
to $g_S=g_0g_1$; so $\gcd(g_S,T_uT_{u'})=1$. Finally $T_u\equiv1\equiv T_{u'}\bmod(g_0g_1)^2$ by (d9), and reducing modulo $g_S$
gives $T_u\equiv T_{u'}\bmod g_S$.
\end{proof}

\begin{remark}\label{pf:sh-misc:standrem}
Of the standing hypotheses, \cref{pf:sh-misc:standof} uses only (a24) together with $C_l>2$; of the standing numerical data
only (d9), (d11), (d12) and (d14). The constant $C_0$ and the size data enter only through these conjuncts.
\end{remark}

\paragraph{The global Theorem B for the mass data.}
We state the result in the form of the formal development; its objects are those of the counting theorems of
\cref{sec:pf:thmS} and \cref{sec:pf:open0}: a datum $D$ with size data $L_2$, a Hida setting $X$ of $D$ with full data $Y$,
a word $u$, and the term system $Y.t_O(u,u)$ (of the diagonal operator of $u$) with its set of generic terms. The
\emph{counting context} $\mathrm{Ctx}(C_g,c_b,C_w,c_T,C_T,C_0,C',C_l;D,L_2,X,Y,u,u)$ is the conjunction of $\mathrm{H}_0(C_0)$ and
of the further counting hypotheses for the pair $(u,u)$ (\cref{sec:pf:thmS}). A \emph{mass datum} for $(X,Y,u)$ consists of an
adelic bridge isomorphism $h_b$, base vectors $X_0$, an additive character $\psi$ of $\A_F$ with the mass-data tie, and the
remaining data $(s_d,C,\mathrm{Th},F_0,\mathcal W,v_R,Q,i_0,i'_0,\nu_q,I_{\mathrm{rat}},I_{\mathrm{re}})$ (a flag $F_0$ at the level prime
$T_u$, a finite set $\mathcal W$ of lattices, a pair system $Q$, measures $\nu_q$ on the archimedean group, reals); it is
\emph{admissible} if the mass data built from it are admissible with the two reals $I_{\mathrm{rat}},I_{\mathrm{re}}$. For a pinned
measure $\nu$ and further measures $\nu_L$ and a function $n_g$ on points, $N=N(Y,u,\nu_L,\nu,n_g)$ denotes the normalised point
data of \cref{sec:pf:thmS}, which attach to every point $\gamma$ (a tuple of $r_S+1$ vectors over $E$) a multiplicity $n(\gamma)$,
two counting factors $N_G(\gamma),N_D(\gamma)$ and a Jacobian bound $J(\gamma)$; $w(\gamma)$ denotes the block-minor weight of $\gamma$.
For $A_{\exp}\in\N$ and $T_B\in\R$ the \emph{point-sum input} $\mathrm{PS}(A_{\exp},T_B)$ for the generic terms of $Y.t_O(u,u)$
and $N$ is the statement: for every finite set $S$ of points each of which occurs in a generic term,
\[
\sum_{\gamma\in S}w(\gamma)\,n(\gamma)\,\bigl(N_G(\gamma)N_D(\gamma)\bigr)^{A_{\exp}}J(\gamma)\ \le\ T_B .
\]

\begin{theorem}[The global Theorem B for the mass data]\label{pf:sh-misc:thmB}
For all reals $C_g,c_b,C_w,c_T,C_T,C_0$ and every $A\in\N$ there are $C_B\ge0$ and reals $C'_0,C_{l,0}$ such that for all
$C'\ge C'_0$ and $C_l\ge C_{l,0}$ the following holds. Let $D$ be a datum with size data $L_2$, $X$ a Hida setting of $D$ with
full data $Y$ and $u$ a word with the counting context $\mathrm{Ctx}(C_g,c_b,C_w,c_T,C_T,C_0,C',C_l;D,L_2,X,Y,u,u)$; let a
mass datum with the mass-data tie be given which is admissible with reals $I_{\mathrm{rat}},I_{\mathrm{re}}$, and let $\nu$ be pinned
for $(X,u,u)$, $\nu_L$ a family of measures and $n_g$ a function on points. Then the point-sum input
\[
\mathrm{PS}\Bigl(A\,(r_S+1)^2,\ \exp\bigl(C_B\,(2r_S+2)^5f\log b\bigr)\,I_{\mathrm{re}}\Bigr)
\]
holds for the generic terms of $Y.t_O(u,u)$ and $N(Y,u,\nu_L,\nu,n_g)$.
\formallabel{Formalism.MainStatements.RowRest.bDat_theoremB_global}
\end{theorem}

The pinned form used in \cref{pf:stageU:mainrow} has in addition a hypothesis on the step choices which is not needed here.

The proof combines four uniform statements:
\begin{enumerate}[label=(\roman*)]
\item\label{pf:sh-misc:B-i} the uniform joint norm bound for rotated points, $\sfname{JointNormBoundRotatedUniform}$: for all $A_w\ge0$ and $C_h$
there are $C_b,C_{g,0}$ such that the rotated form of Theorem~B holds for every datum and every $C_g\ge C_{g,0}$;
\item\label{pf:sh-misc:B-ii} the uniform amplitude statement $\sfname{AmpUnif}$: for all constants there are $C_a\ge0$ and thresholds
beyond which the amplitudes of every datum are Bessel amplitudes with constant $C_a$;
\item\label{pf:sh-misc:B-iii} the uniform point facts $\sfname{PtUnifG}$: for all constants there are thresholds beyond which, under the
counting context, every generic term $t$ of $Y.t_O(u,u)$ with non-zero vector has a \emph{row point} (its point is a point of the row
of $(T_u,T_u)$ at which the local weights are positive);
\item\label{pf:sh-misc:B-iv} the uniform row data $\sfname{RowDataUnifN}$: for all constants there are $C_1,C_n$ and thresholds such that,
under the counting context and for admissible mass data and pinned $\nu$, the row data with constants $(C_1,C_n)$ and $I_{\mathrm{re}}$
exist.
\end{enumerate}

\begin{lemma}[Uniform point facts]\label{pf:sh-misc:ptunif}
Statement \ref{pf:sh-misc:B-iii} holds, with thresholds $C'_0=C_{l,0}=0$.
\formallabel{Formalism.L2.SG.PointUniformity.ptUnifG}
\end{lemma}

\begin{proof}
A generic term $t$ has non-zero weight $G(t)$ (non-vanishing of the weight is one of the conditions defining the generic set), and a
term with $G(t)\ne0$ and non-zero vector has a row point (\lref{Formalism.L2.SG.PointUniformity.rowPt_of}).
\end{proof}

\begin{lemma}[The combination]\label{pf:sh-misc:combo}
Statements \ref{pf:sh-misc:B-i}--\ref{pf:sh-misc:B-iv} imply the assertion of \cref{pf:sh-misc:thmB}, with
$C_B=\max(0,C_b)\max(2,C_w)+|C_a|+|C_n|$ for the constants $C_b,C_a,C_n$ provided by them.
\formallabel{Formalism.MainStatements.RowRest.stepThmBG_ofN}
\end{lemma}

\begin{proof}
Given $C_g,c_b,C_w,c_T,C_T,C_0,A$, let $C_1,C_n$ and thresholds $c_0,d_0$ be given by \ref{pf:sh-misc:B-iv}, $C_a$ and thresholds $a_0,b_0$
by \ref{pf:sh-misc:B-ii} (applied with $C_1$), thresholds $g_0',h_0$ by \ref{pf:sh-misc:B-iii}, and $C_b,C_{g,0}$ by \ref{pf:sh-misc:B-i}
applied with $A_w=A$ and $C_h=|C_T|+\frac18+2|C_g|$. Put
\begin{gather*}
C_B=\max(0,C_b)\max(2,C_w)+|C_a|+|C_n|\ \ (\ge0),\qquad C'_0=\max(a_0,c_0,g_0'),\\
C_{l,0}=\max\bigl(8,128\max(C_{g,0},0),b_0,d_0,h_0\bigr).
\end{gather*}
For $C'\ge C'_0$, $C_l\ge C_{l,0}$ and data as in the statement, the counting context gives $\mathrm{H}_0(C_0)$, the tie
hypotheses for $(u,u)$ and $T_u>0$; with $m\ge C_l\ge128\max(C_{g,0},0)$ and $C_l\ge8$ one obtains the standard hypotheses
in the form required by the core lemma (with $C_s=\max(C_{g,0},0)$ and $C_\beta=\max(2,C_w)$; since $C_\beta\ge2$ and $C_h\ge0$, we have
$C_T\le C_h\le C_\beta C_h$ and $\frac18+2C_g\le C_h\le C_\beta C_h$). Then \ref{pf:sh-misc:B-iv} provides the row data (a Bessel choice
$C_c$ with window and good-gap properties, the measure and point conditions and the norm factorisation), \ref{pf:sh-misc:B-ii} the
Bessel amplitudes, \ref{pf:sh-misc:B-iii} the point tie, and the joint amplitude factorisation for generic terms follows from the row
data (\lref{Formalism.L2.SG.Pf.ThmBChainG.jointAmplitudeRotated_generic}). The core-tie lemma
(\lref{Formalism.MainStatements.RowRest.ties_coreN}) combines these into the ties of the rotated Theorem~B with the bound
$\exp\bigl(C_B(2r_S+2)^5f\log b\bigr)I_{\mathrm{re}}$ (which is $\ge0$ as $I_{\mathrm{re}}\ge0$ by the row data), and the rotated
Theorem~B \ref{pf:sh-misc:B-i}, applied with $C_g=\max(C_{g,0},0)$, turns the ties into the point-sum input
(\lref{Formalism.L2.SG.ThmBUseR.pointSum_of_B3R}).
\end{proof}

\begin{proof}[Proof of \cref{pf:sh-misc:thmB}]
Statement \ref{pf:sh-misc:B-i} is deduced (\lref{Formalism.L2.SG.Pf.ThmBChain.theoremB_uniform_rotated}) from the uniform
weighted bound for rotated points (\lref{CompanionTheory.WaveGrid.weighted_bound_rotated_uniform}); \ref{pf:sh-misc:B-ii} is
\lref{CompanionTheory.Lime.BM.ampUnif_holds}; \ref{pf:sh-misc:B-iii} is \cref{pf:sh-misc:ptunif}; and \ref{pf:sh-misc:B-iv} is deduced
(\lref{Formalism.MainStatements.RowRest.rowDataUnifN_of_6}) from the Bessel bound on the double cosets
(\lref{Formalism.MainNotions.BridgeCS3dBOW.besselOfW_holds}) and the uniform restriction bound for amplitudes
(\lref{Formalism.MainStatements.RowRest.uniform_restriction_bound}). \Cref{pf:sh-misc:combo} concludes.
\end{proof}

\begin{remark}\label{pf:sh-misc:thmBrem}
The deep inputs of \cref{pf:sh-misc:thmB} are \ref{pf:sh-misc:B-i}, \ref{pf:sh-misc:B-ii} and \ref{pf:sh-misc:B-iv}, which are proved
elsewhere; \cref{pf:sh-misc:combo} is a choice of constants.
\end{remark}

\paragraph{Base-level conditions from primitive links.}
The following statement concerns the indefinite real data of the surjectivity argument. Let $A_3$ be a family of the
third level, $R$ a real datum over $A_3$, $F$ a face (subgroups $P_0,P_1$ with characters $\psi_0,\psi_1$, definite column and
period data and a scale), and $R.\mathrm{Cx}(F)$ the complex datum of $R$ at $F$. The \emph{primitive rows} $\mathrm{PR}(R,F)$
are the conjunction of the Hecke and inert rationality conditions for the rows and the lifts, the theta conditions with
their Hecke choices, the unit and equivalence conditions for the elements $g$, the level conditions, the bijectivity of
the principal level map, the lower-level lift type, the block condition, the row and lift tests, positivity and scalar
positivity, the Theorem~A and Theorem~B conditions,
the sum decompositions $E$ and $W$, the uniformity, the lattice condition, the base lift test, the scalar conjugacy, the
dimension count and the surjectivity of the lifts. The \emph{weak characterisation} of $R$ at $F$ contains the model
property at the base, the tuple property and the phase property.

\begin{lemma}[Base-level conditions from primitive links]\label{pf:sh-misc:cdc}
Let $s\in\N$. Assume that $A_3$ satisfies the admissibility condition $\mathrm{Ok}_3(s)$ and the separation of types, and that
$R$ satisfies the weak characterisation and the primitive rows at $F$. Then the complex datum $R.\mathrm{Cx}(F)$ satisfies the
conditions on complex data at the base level, for the data $(P_0,\psi_0,P_1,\psi_1)$, the definite column, period and scale
data of $F$, the integer $s$, the eigensystem choice determined by the left Hecke choice of the theta conditions, the right
and left Hecke choices of the theta conditions, and the actions $g_E$, $g_W$ of $G$ attached to $R$ and $F$.
\formallabel{Formalism.IndefR.Links.complex_data_conditions_base_level_of_primitive_links}
\end{lemma}

\begin{proof}
The conditions at the base level are a list of thirty-seven statements; we indicate where each comes from.
(1) The right and left kernel forms and the right and left $J$-isomorphisms follow from the model property of the weak
characterisation (\lref{Formalism.IndefR.Links.RealDat3.kFormR_of},
\lref{Formalism.IndefR.Links.RealDat3.kFormL_of}, \lref{Formalism.IndefR.Links.RealDat3.jIsoR_of},
\lref{Formalism.IndefR.Links.RealDat3.jIsoL_of}).
(2) The tie of the numbers of points holds by definition; $\mathrm{Ok}_3(s)$ and the separation of types are hypotheses.
(3) The Hecke conditions for rows and lifts, the theta conditions (with the chosen Hecke data), the unit and equivalence
conditions on the right and on the left, the two level conditions, the bijectivity of the principal level map, the
lower-level lift type, the block condition, the row and lift tests, positivity, the sum decompositions $E$ and $W$, the
lattice condition, the base lift test, the Galois condition (scalar conjugacy) and the dimension count are fields of the
primitive rows.
(4) The height change at level zero follows from the tuple property of the weak characterisation
(\lref{Formalism.IndefR.Links.RealDat3.height_change_of_D0}); the positivity of the scalars and of the point numbers from
the scalar positivity of the primitive rows and \lref{Formalism.IndefR.Links.RealDat.nP_pos}.
(5) The Theorem~B operator condition follows from the Theorem~B field of the primitive rows
(\lref{Formalism.IndefR.Links.RealDat3.thmBop_of}); the Theorem~A operator condition from the Theorem~A field and the phase
property (\lref{Formalism.IndefR.Links.RealDat3.thmAop_of}); the kernel form from the uniformity field
(\lref{Formalism.IndefR.Links.RealDat3.kForm_of}); the norm condition $\mathrm{EW}$ holds for every real datum
(\lref{Formalism.IndefR.Links.RealDat3.ewNorm}).
(6) The eigenvalue and type conditions at level zero are the two parts of the unitary-two package
(\lref{Formalism.IndefR.U2.unitary_two_package}), applied to the data above and to the conditions obtained in (1)--(3) (the
kernel forms, the $J$-isomorphisms, the row and lift tests, the left Hecke, theta, unit, equivalence, level-lift and sum
conditions, and the surjectivity of the lifts).
\end{proof}

% ===========================================================================
\subsection{Hosts and their characters}\label{sec:pf:sh-misc:hosts}

This subsection proves five elementary facts about the hosts of the window bound, used in
\cref{sec:pf:vwin,sec:pf:eis,sec:pf:modularity}. We use the notation of \cref{pf:vwin-he:coef,pf:vwin-he:host,pf:vwin-he:cases}
and of \cref{def:window:cases}: a host $h=(\chi,m,q_1,d)$, the conductor $N_\chi$ and order $o_\chi$ of $\chi$,
$N_h=\lcm(2^{m+2},N_\chi)q_1d$, the subgroup $H_h\subset(\Z/N_h)^\times$, the primes $p_K=p_K(\chi)$ and $p_1=p_1(\chi)$, the class number
$h^*=h^*(\chi)$, the bad set $B_\chi=\{3\}\cup\{\ell\text{ odd prime}:\ell\mid P_\chi\}$ with $P_\chi=o_\chi N_\chi p_Kh^*p_1$, and, for a Dirichlet
character $\psi$ modulo $N_h$ with values in a commutative ring, its components $\psi_{(r)}$, its $\ell$-power-order part
$\psi^{(\ell)}$ and prime-to-$\ell$ part $\psi_{\ell'}$ (so $\psi=\psi_{\ell'}\psi^{(\ell)}$), $\psi^\flat=\psi_{\ell'}(\psi_{\ell'})_{(\ell)}^{-1}$ and
$c_2(\psi)=v_2(\cond\psi_{(2)})$. The cases I, II, III, V of $(\ell,j,\psi)$ are those of \cref{pf:vwin-he:cases}; following
\cref{tab:window:cases}, $(\ell,j,\psi)$ is in case Z if $\ell\notin B_\chi$, $(\ell-1)\mid2j$ and $\psi=1$, and in case IV if
$\ell\in B_\chi$, $j\ge2$ and $\ord\psi^\flat$ is odd. For an integer $n$ and a prime $p$ we write $n_p$ for the $p$-part
$p^{v_p(n)}$ of $n$ and $n_{(2')}$ for the odd part of $n$.

We use the following shared facts, proved elsewhere.
\begin{enumerate}[label=(F\arabic*)]
\item\label{pf:sh-misc:F-comp} (\lref{UpperTheory.V.HV.compAlg}, \lref{UpperTheory.V.HV.compCond},
\lref{UpperTheory.V.HV.compOrd}) For characters modulo $N$ with values in a domain and a prime $r$: $\psi\mapsto\psi_{(r)}$ is a group
homomorphism with $\ord\psi_{(r)}\mid\ord\psi$; $\psi=\prod_{p\mid N}\psi_{(p)}$; $\cond\psi_{(r)}=(\cond\psi)_r$, so $\psi_{(r)}=1$ if and only
if $r\nmid\cond\psi$; for $M\mid N$, $\cond\psi\mid M$ if and only if $\cond\psi_{(p)}\mid M$ for every prime $p\mid N$; the order of
$\psi_{(2)}$ is a power of $2$.
\item\label{pf:sh-misc:F-part} (\lref{UpperTheory.V.HV.partOrd}) For an element $x$ of finite order of a group and a prime $\ell$,
$\ord x^{(\ell)}=(\ord x)_\ell$ and $\ord x_{\ell'}=\ord x/(\ord x)_\ell$; if $\ell\nmid\ord x$, then $x^{(\ell)}=1$ and $x_{\ell'}=x$.
\item\label{pf:sh-misc:F-hostcomp} (\lref{UpperTheory.V.HV.hostComp}) If $H_h\subseteq\ker\psi$, then $\ord\psi\mid2^{m+1}o_\chi$, and for
every odd prime $r$: $\ord\psi_{(r)}\mid2o_\chi$; if $r\nmid N_\chi$ then $\psi_{(r)}^2=1$; if $r\nmid N_h$ then $\psi_{(r)}=1$.
\item\label{pf:sh-misc:F-th1} (\lref{UpperTheory.V.HV.hostTh1}) If $H_h\subseteq\ker\psi$ and $\ell$ is an odd prime, then
$(\psi^\flat)_{(2)}=\psi_{(2)}$ and $(\psi^\flat)_{(r)}=(\psi_{\ell'})_{(r)}$ for every prime $r\ne\ell$.
\item\label{pf:sh-misc:F-pK} (\lref{UpperTheory.V.pK_pOne_prime_conditions}) $p_K$ is a prime $\equiv3\bmod4$, $p_K\ge7$,
$p_K\nmid N_\chi$; $p_1$ is a prime, $p_1\equiv1$ modulo the odd part of $o_\chi$, and $p_1\nmid2N_\chi p_K$.
\end{enumerate}

\begin{lemma}[The level of a host]\label{pf:sh-misc:hostindex}
Let $h=(\chi,m,q_1,d)$ be a host. Then:
\begin{enumerate}[label=(\alph*)]
\item the index $[(\Z/N_h)^\times:H_h]$ divides $2^{m+2}o_\chi$ and is at most $N_h$;
\item $16\mid N_h$, and $0<o_\chi\le N_\chi$;
\item a prime $p$ divides $N_h$ if and only if $p=2$, $p\mid N_\chi$, $p=q_1$ or $p=d$;
\item $(N_h)_p=(N_\chi)_p$ for every prime $p\notin\{2,q_1,d\}$, and $(N_h)_{q_1}=q_1$, $(N_h)_d=d$;
\item $u^{2^{m+1}o_\chi}\in H_h$ for every $u\in(\Z/N_h)^\times$.
\end{enumerate}
\formallabel{UpperTheory.V.HV.hostIndex}
\end{lemma}

\begin{proof}
Recall that $q_1,d$ are primes with $q_1\equiv3\bmod8$, $d\equiv7\bmod16$ and $q_1d$ prime to $N_\chi$. Hence $q_1,d\ne2$, $q_1\ne d$
(they differ modulo $8$), and $q_1,d\nmid N_\chi$. For every prime $p$,
\[
v_p(N_h)=\max\bigl((m+2)\,[p=2],\,v_p(N_\chi)\bigr)+[p=q_1]+[p=d],
\]
where $[\cdot]$ is $1$ if the condition holds and $0$ otherwise. This gives (c) and (d) (for $p\notin\{2,q_1,d\}$ the formula gives
$v_p(N_\chi)$, and $v_{q_1}(N_\chi)=v_d(N_\chi)=0$). As $m\ge2$, $16\mid2^{m+2}\mid N_h$. The order $o_\chi$ of $\chi$ is positive and at most
$\varphi(N_\chi)\le N_\chi$ (the values of $\chi$ on the units are $\varphi(N_\chi)$-th roots of unity), which proves (b).

(a) $H_h$ is the intersection of the preimages, under the reduction maps from $(\Z/N_h)^\times$, of $\{\pm1\}\subset(\Z/2^{m+2})^\times$, of
$\ker\chi\subset(\Z/N_\chi)^\times$, and of the subgroups of squares in $(\Z/q_1)^\times$ and $(\Z/d)^\times$. In a commutative group the index
of a preimage $f^{-1}(K)$ under a homomorphism divides the index of $K$, and the index of an intersection of two subgroups divides the
product of their indices. The four indices divide $2^m$ (the group $(\Z/2^{m+2})^\times$ has order $2^{m+1}$), $o_\chi$ (the quotient by
$\ker\chi$ is the image of $\chi$, a group of roots of unity of order dividing $o_\chi$), $2$ and $2$. Hence the index divides
$2^m\cdot o_\chi\cdot2\cdot2=2^{m+2}o_\chi$. It also divides $|(\Z/N_h)^\times|=\varphi(N_h)\le N_h$.

(e) Let $u\in(\Z/N_h)^\times$ and $w=u^{2^{m+1}o_\chi}$. Then $w=(u^{2^mo_\chi})^2$ is a square, so its reductions modulo $q_1$ and $d$ are
squares; its reduction modulo $2^{m+2}$ is $1\in\{\pm1\}$ because $\varphi(2^{m+2})=2^{m+1}$ divides the exponent; and
$\chi(w)=(\chi(u)^{o_\chi})^{2^{m+1}}=1$. Hence $w\in H_h$.
\end{proof}

\begin{lemma}[Characters of odd order]\label{pf:sh-misc:hostodd}
Let $h=(\chi,m,q_1,d)$ be a host, $O$ a domain and $\psi$ a character modulo $N_h$ with values in $O$ and $H_h\subseteq\ker\psi$.
\begin{enumerate}[label=(\alph*)]
\item If $\ord\psi$ is odd, then $\ord\psi$ divides the odd part of $o_\chi$, $\psi_{(2)}=1$, $c_2(\psi)=0$, and $\cond\psi$ is odd and
divides $N_\chi$.
\item If $\ell$ is a prime with $\ell\ne2,q_1,d$ and $\ord\psi^\flat$ is odd, then $\cond\psi\mid N_\chi$ and $\psi_{(2)}=1$.
\end{enumerate}
\formallabel{UpperTheory.V.HV.hostOdd}
\end{lemma}

\begin{proof}
We first show: \emph{(C) if $\psi_{(2)}=1$ and, for $p\in\{q_1,d\}$, the order of $\psi_{(p)}$ divides an odd integer, then $\cond\psi\mid
N_\chi$.} By \ref{pf:sh-misc:F-comp} it suffices to show $\cond\psi_{(p)}\mid N_\chi$ for every prime $p\mid N_h$. For $p=2$,
$\psi_{(2)}=1$ has conductor $1$. For $p\in\{q_1,d\}$, $\cond\psi_{(p)}=(\cond\psi)_p$ divides $(N_h)_p=p$
(\cref{pf:sh-misc:hostindex}(d)); if $p\mid N_\chi$ we are done, and otherwise $\psi_{(p)}^2=1$ by \ref{pf:sh-misc:F-hostcomp},
so $\ord\psi_{(p)}$ divides $2$ and an odd integer, hence $\psi_{(p)}=1$. For the remaining primes $p\mid N_h$, $p\mid N_\chi$ by
\cref{pf:sh-misc:hostindex}(c), and $\cond\psi_{(p)}$ divides $(N_h)_p=(N_\chi)_p$ (\cref{pf:sh-misc:hostindex}(d)), which divides
$N_\chi$.

Also, \emph{if $\ord x$ is odd for a character $x$, then $x_{(2)}=1$}: the order of $x_{(2)}$ is a power of $2$ and divides $\ord x$
(\ref{pf:sh-misc:F-comp}).

(a) By the last remark $\psi_{(2)}=1$. For $p\in\{q_1,d\}$, $\ord\psi_{(p)}$ divides the odd integer $\ord\psi$; so (C) gives
$\cond\psi\mid N_\chi$. By \ref{pf:sh-misc:F-comp}, $\psi_{(2)}=1$ means $2\nmid\cond\psi$, so $\cond\psi$ is odd and
$c_2(\psi)=v_2(\cond\psi_{(2)})=v_2((\cond\psi)_2)=0$. Finally $\ord\psi\mid2^{m+1}o_\chi$ by \ref{pf:sh-misc:F-hostcomp}; as
$\ord\psi$ is odd it divides $o_\chi$, hence the odd part of $o_\chi$.

(b) By \ref{pf:sh-misc:F-th1}, $\psi_{(2)}=(\psi^\flat)_{(2)}$, which is trivial because $\ord\psi^\flat$ is odd. For $p\in\{q_1,d\}$ we
have $p\ne\ell$, and $\psi_{(p)}=(\psi_{\ell'})_{(p)}(\psi^{(\ell)})_{(p)}=(\psi^\flat)_{(p)}(\psi^{(\ell)})_{(p)}$ by
\ref{pf:sh-misc:F-comp} and \ref{pf:sh-misc:F-th1}. The first factor has order dividing $\ord\psi^\flat$ and the second order dividing
$\ord\psi^{(\ell)}$, a power of the odd prime $\ell$ (\ref{pf:sh-misc:F-part}); in a commutative group the order of a product divides the
product of the orders, so $\ord\psi_{(p)}$ divides the odd integer $\ord\psi^\flat\cdot\ell^{v_\ell(\ord\psi)}$. Now (C) gives
$\cond\psi\mid N_\chi$.
\end{proof}

\begin{lemma}[Bad primes of big hosts]\label{pf:sh-misc:bchi}
Let $h=(\chi,m,q_1,d)$ be a host which is big for $t\in\N$. Then:
\begin{enumerate}[label=(\alph*)]
\item every $\ell\in B_\chi$ satisfies $\ell<q_1$ and $\ell<d$;
\item $p_K\in B_\chi$, $p_K\nmid N_h$ and $p_1\nmid N_h$;
\item every odd prime $\ell\notin B_\chi$ satisfies $\ell\ge5$, $\ell\nmid o_\chi$, $\ell\nmid N_\chi$, $\ell\ne p_K$, $\ell\nmid h^*$, $\ell\ne p_1$ and
$\ell\nmid[(\Z/N_h)^\times:H_h]$; and for every $0<j<t$ with $(\ell-1)\mid2j$ one has $j\ge2$, $\ell\le2t$ and $\ell\nmid N_h$.
\end{enumerate}
\formallabel{UpperTheory.V.HV.bchiFacts}
\end{lemma}

\begin{proof}
Bigness means $2t+1<q_1$, $2t+1<d$, and $N_\chi,o_\chi,p_K,h^*,p_1<\min(q_1,d)$. The integer $P_\chi=o_\chi N_\chi p_Kh^*p_1$ is
non-zero (class numbers are positive), so every odd prime dividing $P_\chi$ lies in $B_\chi$, and every prime divisor of one of the
five factors is at most that factor.

(a) If $\ell=3$, then $3<7\le p_K<\min(q_1,d)$ by \ref{pf:sh-misc:F-pK}. Otherwise $\ell$ is a prime dividing $P_\chi$, so it divides
one of the five factors and is at most that factor, which is smaller than $\min(q_1,d)$.

(b) $p_K$ is an odd prime dividing $P_\chi$, so $p_K\in B_\chi$. By \cref{pf:sh-misc:hostindex}(c) a prime divisor of $N_h$ is $2$, a
divisor of $N_\chi$, $q_1$ or $d$; but $p_K\ge7$, $p_K\nmid N_\chi$ and $p_K<\min(q_1,d)$, so $p_K\nmid N_h$. Likewise $p_1\nmid2N_\chi$
and $p_1<\min(q_1,d)$, so $p_1\nmid N_h$.

(c) Since $3\in B_\chi$, $\ell\ne3$, so $\ell\ge5$. If $\ell$ divided one of $o_\chi,N_\chi,p_K,h^*,p_1$, it would divide $P_\chi$ and lie in
$B_\chi$; this gives the five non-divisibility statements (for the primes $p_K,p_1$, $\ell\mid p$ means $\ell=p$). The index of $H_h$
divides $2^{m+2}o_\chi$ (\cref{pf:sh-misc:hostindex}(a)), and $\ell$ is odd with $\ell\nmid o_\chi$, so $\ell$ does not divide the index.
Let $0<j<t$ with $(\ell-1)\mid2j$. Then $4\le\ell-1\le2j$, so $j\ge2$, and $\ell\le2j+1\le2t-1\le2t$. Finally $\ell\ne2$,
$\ell\nmid N_\chi$, and $\ell\le2t<\min(q_1,d)$, so $\ell\nmid N_h$ by \cref{pf:sh-misc:hostindex}(c).
\end{proof}

\begin{proposition}[The case division]\label{pf:sh-misc:cases}
Let $h$ be a host which is big for $t$, $O$ a domain, $\psi$ a character modulo $N_h$ with values in $O$ and $H_h\subseteq\ker\psi$,
$\ell$ an odd prime and $0<j<t$.
\begin{enumerate}[label=(\alph*)]
\item If $\ell\notin B_\chi$, then $\ell\nmid\ord\psi$, $\psi_{(\ell)}^2=1$, $\psi_{\ell'}=\psi$, $\psi^{(\ell)}=1$; if moreover $(\ell-1)\mid2j$, then
$\psi^\flat=\psi$; and $(\ell,j,\psi)$ is in case I, Z, II or III.
\item If $\ell\in B_\chi$, then $(\ell,j,\psi)$ is in case II, IV or V.
\end{enumerate}
\formallabel{UpperTheory.V.HV.cases}
\end{proposition}

\begin{proof}
(a) By \cref{pf:sh-misc:bchi}(c), $\ell\nmid o_\chi$ and $\ell\nmid N_\chi$. By \ref{pf:sh-misc:F-hostcomp}, $\ord\psi\mid2^{m+1}o_\chi$, and
$\ell$ is odd, so $\ell\nmid\ord\psi$; by \ref{pf:sh-misc:F-part}, $\psi^{(\ell)}=1$ and $\psi_{\ell'}=\psi$. Since $\ell\nmid N_\chi$,
$\psi_{(\ell)}^2=1$ by \ref{pf:sh-misc:F-hostcomp}. If $(\ell-1)\mid2j$, then $\ell\nmid N_h$ by \cref{pf:sh-misc:bchi}(c), so
$\psi_{(\ell)}=1$ by \ref{pf:sh-misc:F-hostcomp}, and $\psi^\flat=\psi_{\ell'}(\psi_{\ell'})_{(\ell)}^{-1}=\psi\,\psi_{(\ell)}^{-1}=\psi$.
If $(\ell-1)\nmid2j$, the triple is in case I. Let $(\ell-1)\mid2j$. If $\psi=1$, it is in case Z; if $\psi\ne1$ and $\ord\psi$ is odd, in case
III; if $\ord\psi$ is even, then $j\ge2$ (\cref{pf:sh-misc:bchi}(c)) and $\ord\psi^\flat=\ord\psi$ is even, so it is in case II.

(b) If $j=1$, the triple is in case V. Otherwise $j\ge2$, and it is in case II if $\ord\psi^\flat$ is even and in case IV if it is odd.
\end{proof}

\begin{lemma}[Coefficient rings]\label{pf:sh-misc:coefsup}
Let $\ell$ be a prime.
\begin{enumerate}[label=(\alph*)]
\item Any two coefficient rings for $\ell$ are contained in a third.
\item For every $n\ge1$ with $\ell\nmid n$ there is a coefficient ring $\OO$ with $e_\OO=1$ which contains a primitive $n$-th root of unity
(so that the $n$-th roots of unity in $\OO$ form a cyclic group of order $n$).
\item For every coefficient ring $\OO$ and every $x\in\Qlbar$ integral over $\Zl$ there is a coefficient ring $\OO'\supseteq\OO$ with
$x\in\OO'$.
\item If $\OO$ is a coefficient ring, $x\in\Qlbar$ is integral over $\Zl$, and $xb=a$ for some $a,b\in\OO$ with $b\ne0$, then $x\in\OO$.
\end{enumerate}
\formallabel{UpperTheory.V.HV.coefSup}
\end{lemma}

\begin{proof}
\emph{Step 1: every $\Zl$-subalgebra $A\subset\Qlbar$ which is a finitely generated $\Zl$-module is contained in a coefficient ring.}
Let $s$ be a finite set of $\Zl$-module generators of $A$ and $K=\Ql(s)\subset\Qlbar$, a finite extension of $\Ql$ (the elements of
$\Qlbar$ are algebraic over $\Ql$). The integral closure $\OO_K$ of $\Zl$ in $K$ is a discrete valuation ring, finitely generated as a
$\Zl$-module (\cite[Ch.~II, \S2]{Serre1979}), hence a coefficient ring. Every element of $A$ is integral over
$\Zl$ (as $A$ is a finite $\Zl$-module) and lies in $K$, so $A\subseteq\OO_K$.

(a) and (c): the $\Zl$-algebras $\OO\OO'$ and $\OO[x]$ are finitely generated $\Zl$-modules (sums and products of finitely many finite
modules, and $x$ integral); apply Step 1.

(b) Let $\zeta\in\Qlbar$ be a primitive $n$-th root of unity and $A=\Zl[\zeta]$, a finite $\Zl$-algebra. Since $n$ is a unit in $\Zl$,
$A$ is unramified over $\Zl$: in the module of K\"ahler differentials, $0=d(\zeta^n)=n\zeta^{n-1}d\zeta$ with $n\zeta^{n-1}$ a unit, so
$d\zeta=0$, and $d$ vanishes on $A$, whence $\Omega_{A/\Zl}=0$. By Step 1, $A$ is contained in a coefficient ring $\OO_1$; since
$\OO_1$ is integral over $A$, every maximal ideal of $A$ is the contraction of a maximal ideal of $\OO_1$ (lying over), so $A$ is
local, and formal unramifiedness gives that its maximal ideal is generated by $\ell$. A noetherian local domain which is not a field and whose maximal ideal is principal is a discrete
valuation ring; so $A$ is a coefficient ring, and $e_A=\len_A(A/\ell A)=1$. It contains $\zeta$.

(d) The element $a/b$ of the fraction field of $\OO$ maps to $x$ under the embedding of the fraction field into $\Qlbar$; it is
integral over $\Zl\subseteq\OO$, hence over $\OO$. A discrete valuation ring is integrally closed, so $a/b\in\OO$, i.e.\ $x\in\OO$.
\end{proof}

% ===========================================================================
\subsection{Host frames over \texorpdfstring{$\Q(i)$}{Q(i)}}\label{sec:pf:sh-misc:frames}

This subsection proves two facts used in case V of the window bound (\cref{sec:pf:eis:countV}) and in
\cref{sec:pf:modularity}: the Galois group of $\Q(i)$ is the kernel of the quadratic character $\eta_4$, and host frames satisfy the
condition $(\mathrm{F1})$ of \cref{pf:eis:framecond}. Throughout, $\Omega$ is an algebraic closure of $\Q$ and $G_\Q=\Aut(\Omega)$ with the
Krull topology; for $N\ge1$, $\kappa_N\colon G_\Q\to(\Z/N)^\times$ is the cyclotomic character modulo $N$ ($g\zeta=\zeta^{\kappa_N(g)}$ for
$\zeta^N=1$), and for a Dirichlet character $\chi$ modulo $N$ we put $\chi_G=\chi\circ\kappa_N$. The character $\eta_4$ is $(\chi_4)_G$ for the
non-trivial character $\chi_4$ modulo $4$. A number field $K$ is \emph{of type $\Q(i)$} if $[K:\Q]=2$ and $K$ contains an element $i$ with
$i^2=-1$.

\begin{proposition}[The Gaussian field]\label{pf:sh-misc:gauss}
There is a CM field of type $\Q(i)$. For every number field $K\subset\Omega$ of type $\Q(i)$, the image of $G_K=\Aut_K(\Omega)$ in $G_\Q$ is
$\ker\eta_4$, and the discriminant of $K$ is $-4$.
\formallabel{UpperTheory.V.UnitaryThreeInertia.gaussian_field_range_eq_ker_quartic_char}
\end{proposition}

\begin{proof}
\emph{Existence.} The cyclotomic field $\Q(\zeta_4)$ has degree $\varphi(4)=2$, is a CM field (cyclotomic fields other than $\Q$ are
CM), and $\zeta_4^2=-1$: since $\zeta_4^4=1$ and $\zeta_4^2\ne1$, $(\zeta_4^2-1)(\zeta_4^2+1)=0$ gives $\zeta_4^2=-1$.

\emph{Galois group.} Let $K\subset\Omega$ be of type $\Q(i)$, with $i\in K$, $i^2=-1$. Since $i\notin\Q$ and $[K:\Q]=2$, $K=\Q(i)$. The image
$\zeta$ of $i$ in $\Omega$ satisfies $\zeta^2=-1$, so it is a primitive fourth root of unity, and $\eta_4(\sigma)=\chi_4(\kappa_4(\sigma))=1$
if and only if $\kappa_4(\sigma)=1$, i.e.\ if and only if $\sigma(\zeta)=\zeta$. As $K$ is generated by $i$, an automorphism $\sigma$ of $\Omega$
fixes $K$ pointwise if and only if it fixes $\zeta$. Hence the image of $G_K$ is $\ker\eta_4$.

\emph{Discriminant.} $K=\Q(\zeta_4)$ is the cyclotomic field of conductor $2^2$, whose discriminant is
$(-1)^{\varphi(4)/2}2^{2^{1}(2\cdot2-2-1)}=-4$ by the discriminant formula for prime-power cyclotomic fields
\cite[Ch.~2]{Washington1997}.
\end{proof}

\begin{theorem}[Host frames satisfy (F1)]\label{pf:sh-misc:hostframe}
Let $\ell$ be a prime, $\OO$ a coefficient ring for $\ell$, $N\ge1$, $e\in\N$, $\psi$ a Dirichlet character modulo $N$ with values in $\OO$ and
$o_1\in\N$ such that $\mathrm{H}(\ell,e,\psi,o_1)$ holds. Then every host frame $F$ for these data (\cref{pf:eis:frame}(c)) satisfies
$(\mathrm{F1})$.
\formallabel{UpperTheory.V.UnitaryThreeInertia.hostFrame_ok}
\end{theorem}

We use the following facts. (G1) For every $N\ge1$, $\ker\kappa_N$ is open; $\kappa_N(\sigma)=-1$ for every complex conjugation $\sigma$ (an
automorphism with $j\circ\sigma=\overline{j}$ for some embedding $j\colon\Omega\to\C$); and $\kappa_N$ is trivial on the inertia group $I_A$ of every
place $A$ of $\Omega$ not dividing $N$ (\lref{UpperTheory.V.Chr.clsFacts}). (G2) Integers prime to $\ell$ are units in $\OO$, and
$\ell$ lies in the maximal ideal of $\OO$ (\lref{UpperTheory.V.Base.dvrFacts}). (G3) Decomposition and inertia groups of places of
$\Omega$ are closed in $G_\Q$, $I_A$ is normal in $D_A$, and the projection $G_\Q\to G_{\Q,S}$ is continuous and a quotient map
(\lref{ReciprocityTheory.Gal.isClosed_Dec}, \lref{ReciprocityTheory.Gal.isClosed_Ine}). (G4) The $\ell$-adic cyclotomic
character $\chi_{\mathrm{cyc}}$ is trivial on $I_A$ for every place $A$ not above $\ell$. (G5) For every place $A$ above a prime $p$ there is a
Frobenius element $\gamma\in D_A$ ($\gamma(x)\equiv x^p$ modulo the maximal ideal of $A$ for $x\in A$), and every open subgroup of $G_\Q$ which
contains $I_A$ and $\gamma$ contains $D_A$ (the image of $\gamma$ topologically generates $D_A/I_A\cong\Gal(\overline{\F}_p/\F_p)$)
(\cite[Ch.~I, \S9]{Neukirch1999}, passing to the limit over the finite Galois subextensions).

\begin{proof}
Let $S$, places $A_v$ ($v\in S$), $\pi$, $\eta$, $\Psi$, $c$, $D_v=D_{A_v}$, $I_v=I_{A_v}$, $v_\ell$ and $S_b$ be as in the definition of a host frame; so
$\eta\circ\pi=\eta_4$, $\Psi\circ\pi=\Psi_{\psi,e}=\chi_{\mathrm{cyc}}^e\psi_G$ (with values in $\OO^\times$), $c=\pi(\sigma)$ for a complex
conjugation $\sigma$, $v_\ell$ is the place above $\ell$, and $v\in S_b$ if and only if $v$ lies over a prime of
$S_b(\ell,\psi)=\{p\mid N:\ p\ne\ell,\ (\psi^{(\ell)})_{(p)}\ne1\}$. We verify the conditions of $(\mathrm{F1})$.

\emph{$2\in\OO^\times$} by (G2), as $\ell\ne2$.

\emph{Continuity.} $\pi$ is continuous by (G3). $\eta_4=(\chi_4)_G$ has open kernel by (G1), hence is locally constant; since $\eta\circ\pi=\eta_4$
and $\pi$ is a quotient map, $\eta$ is locally constant. For $n\ge0$, $g\mapsto\Psi_{\psi,e}(g)\bmod\mathfrak m_\OO^n$ is locally constant: its
kernel contains the open subgroup of the $g$ with $\psi_G(g)=1$ and $|\chi_{\mathrm{cyc}}(g)-1|_\ell<\ell^{-n+1}$, because for such $g$ the
element $\chi_{\mathrm{cyc}}(g)-1$ lies in $\ell^n\Zl$, hence its image lies in $\ell^n\OO\subseteq\mathfrak m_\OO^n$ (G2). Again this descends
to $\Gamma=G_{\Q,S}$.

\emph{Complex conjugation.} $\sigma^2=1$: if $j\circ\sigma=\bar j$, then $j(\sigma^2x)=\overline{\overline{j(x)}}=j(x)$ and $j$ is injective.
Hence $c^2=\pi(\sigma^2)=1$. Next $\eta(c)=\eta_4(\sigma)=\chi_4(\kappa_4(\sigma))=\chi_4(-1)=-1$ by (G1). Finally
$\chi_{\mathrm{cyc}}(\sigma)=-1$ (by (G1) for the powers of $\ell$), and $\psi_G(\sigma)=\psi(-1)=(-1)^e$ by the parity condition of
$\mathrm{H}(\ell,e,\psi,o_1)$; so $\Psi(c)=(-1)^e(-1)^e=1$.

\emph{$v_\ell\notin S_b$.} Otherwise the prime ideal $v_\ell$ of $\Z$ would contain $\ell$ and a prime $p\in S_b(\ell,\psi)$, $p\ne\ell$; but two
coprime integers cannot lie in the same prime ideal.

\emph{Decomposition and inertia groups.} $D_v$ and $I_v$ are closed, and $I_v$ is a normal subgroup of $D_v$, by (G3).

\emph{Frobenius generators.} Let $v\in S$, $p$ the prime below $v$, and $\gamma\in D_v$ a Frobenius element (G5). The closure $C$ of the subgroup
generated by $I_v$ and $\gamma$ is the intersection of the open subgroups of the profinite group $G_\Q$ containing it, each of which contains
$D_v$ by (G5); so $C\supseteq D_v$, i.e.\ $I_v$ and $\gamma$ topologically generate $D_v$. Let $n\ge1$ and $M=p^n-1$. Then $p\nmid M$, so
$\kappa_M$ is trivial on $I_v$ (G1). The $M$-th roots of unity are distinct modulo the maximal ideal of $A_v$ ($M$ is prime to $p$), so
$\gamma(\zeta)=\zeta^p$ for them, i.e.\ $\kappa_M(\gamma)=p$, an element of order exactly $n$ in $(\Z/M)^\times$ (since $p^k-1<M$ for
$0<k<n$). The preimage under $\kappa_M$ of the cyclic subgroup $\langle p\rangle$ is open and contains $I_v$ and $\gamma$, hence $D_v$
(G5). So $\kappa_M$ maps $D_v$ to $\langle p\rangle\cong\Z/n$, with $p\mapsto1$; the composite $d\colon D_v\to\Z/n$ is locally constant (its
kernel contains $D_v\cap\ker\kappa_M$), trivial on $I_v$, and $d(\gamma)=1$.

\emph{$\eta$ on inertia at $T_b$.} Let $v\in T_b=\{v_\ell\}\cup S_b$. Then $v$ lies over $\ell$ or over a prime of $S_b(\ell,\psi)$; in both cases
over an odd prime (a prime $p=2$ in $S_b(\ell,\psi)$ would require $(\psi^{(\ell)})_{(2)}\ne1$, but this character has order a power of
$\ell$ and a power of $2$ (\ref{pf:sh-misc:F-comp}), so it is trivial). Thus $4$ is a unit at $A_v$, $\kappa_4$ is trivial on $I_v$ (G1),
and so is $\eta_4$.

\emph{The elements $x_v$.} For $v\in T_b$ there is $x\in D_v$ with $\eta(x)=1$ and $\Psi(x)^2-1\in\OO^\times$; this is
\lref{UpperTheory.V.UnitaryThreeInertia.exists_sigma} (which states that for every host frame and $v\in T_b$ there is $x\in D_v$ with
$\eta(x)=1$, $\Psi(x)^2-1\in\OO^\times$ and $\chi_{\mathrm{cyc}}(x)\equiv1\bmod\ell$).

\emph{$\Psi$ on $N_v$ for $v\notin T_b$.} Let $v\in S\setminus T_b$ and $s\in N_v$, i.e.\ $s\in I_v$ with $\eta(s)=1$ and $\Psi(s)\equiv1$ modulo
$\mathfrak m_\OO$. As $v\in S$, the prime $p$ below $v$ divides $\ell N$; $p\ne\ell$ because $v\ne v_\ell$ and there is only one place above $\ell$;
so $p\mid N$, and $(\psi^{(\ell)})_{(p)}=1$ since $v\notin S_b$. By (G4), $\chi_{\mathrm{cyc}}(s)=1$, so $\Psi(s)=\psi_G(s)$. For $s$ in the inertia
group at $p$, $\kappa_N(s)$ is congruent to $1$ modulo the prime-to-$p$ part of $N$, so $\chi_G(s)=(\chi_{(p)})_G(s)$ for every character $\chi$
modulo $N$; in particular $(\psi^{(\ell)})_G(s)=1$ and $\Psi(s)=(\psi_{\ell'})_G(s)=:u$. Then $u^n=1$ for $n=\ord\psi_{\ell'}$, which is prime to
$\ell$ (\ref{pf:sh-misc:F-part}), hence a unit in $\OO$ (G2), and $u\equiv1\bmod\mathfrak m_\OO$. This forces $u=1$: otherwise
$(u-1)\sum_{i<n}u^i=u^n-1=0$ gives $\sum_{i<n}u^i=0$, whose reduction is $n\equiv0\bmod\mathfrak m_\OO$, contradicting that $n$ is a unit.
\end{proof}

% ===========================================================================
\subsection{Two facts on smooth representations of \texorpdfstring{$\GL_n$}{GLn}}\label{sec:pf:sh-misc:llc}

This subsection proves the facts \ref{pf:llc:F-twtame} and \ref{pf:llc:F-goodseg} used in the construction of local parameter maps
(\cref{sec:pf:llc}) and in the Galois reduction (\cref{sec:pf:galred}). We use the notation of \cref{pf:llc:sec-setting}, for an
arbitrary non-archimedean local field $k$ of characteristic zero with ring of integers $\OO_k$, maximal ideal $\mathfrak p_k$ and
uniformiser $\varpi$: $\shmiscIrr(k)$, $n(y)$, $\pi_y$, the twist $y\otimes\chi$ by a smooth character $\chi$ of $k^\times$, $\nu^s=|\cdot|^s$,
$\shmiscKone=\ker(\GL_n(\OO_k)\to\GL_n(\OO_k/\mathfrak p_k))$, depth zero (a non-zero $\shmiscKone$-fixed vector), tame characters ($\chi(u)=1$ for
$u\in1+\mathfrak p_k$), segments $\Delta$ with cuspidal $\rho_\Delta$ and length $\ell_\Delta$, $x=L(l)$, supercuspidal supports, condition
\shmiscGo{} and good representations (\cref{pf:llc:def-good}).

\begin{proposition}[Depth zero and tame twists]\label{pf:sh-misc:twtame}
Let $y\in\shmiscIrr(k)$ with $n(y)>0$ and of depth zero, and let $\chi$ be a smooth character of $k^\times$. Then $y\otimes\chi$ has depth zero
if and only if $\chi$ is tame.
\formallabel{UpperTheory.Zero.LLC.twTame}
\end{proposition}

\begin{proof}
Write $n=n(y)$, $\pi=\pi_y$ on $V$, and $\sigma=\pi\otimes(\chi\circ\det)$, a representative of $y\otimes\chi$. We use two observations.
(i) If $g\in\shmiscKone$, then $\det g\in1+\mathfrak p_k$: the entries of $g$ are integral and $g\equiv1$ modulo $\mathfrak p_k$, so
$\det g\equiv\det1=1$. (ii) For $u\in1+\mathfrak p_k$ let $d_u=\mathrm{diag}(u,1,\dots,1)$ (here $n>0$ is used); then $d_u\in\shmiscKone$,
$\det d_u=u$, and $d_u$ commutes with every diagonal matrix.

\emph{If $\chi$ is tame}, then by (i) $\chi(\det g)=1$ for $g\in\shmiscKone$, so the $\shmiscKone$-fixed vectors of $\pi$ and of $\sigma$ coincide, and a non-zero
$\shmiscKone$-fixed vector of $\pi$ is one of $\sigma$; isomorphic representations have non-zero fixed vectors simultaneously.

\emph{Conversely}, assume that $y\otimes\chi$ has depth zero and that $c=\chi(u)\ne1$ for some $u\in1+\mathfrak p_k$. Let
$V(\shmiscKone)\subseteq V$ be the span of the vectors $\pi(k)z-z$, $k\in\shmiscKone$, $z\in V$. Let $0\ne v\in V^{\shmiscKone}$ (depth zero of $y$) and
$0\ne w\in V$ fixed by $\sigma(\shmiscKone)$.

\emph{Step 1: if $x\in V$ is fixed by $\sigma(\shmiscKone)$, then $\pi(t)x\in V(\shmiscKone)$ for every diagonal matrix $t=\mathrm{diag}(\varpi^{\mu_1},
\dots,\varpi^{\mu_n})$.} Indeed $\chi(u)\pi(d_u)x=\sigma(d_u)x=x$, so $\pi(d_u)x=c^{-1}x$, and since $d_u$ commutes with $t$,
$\pi(d_u)\pi(t)x=c^{-1}\pi(t)x$. Hence $(c^{-1}-1)\pi(t)x=\pi(d_u)\pi(t)x-\pi(t)x\in V(\shmiscKone)$, and $c^{-1}\ne1$.

\emph{Step 2: $\pi(g)w\in V(\shmiscKone)$ for every $g\in\GL_n(k)$.} By the Cartan decomposition, $g=k_1tk_2$ with $k_1,k_2\in\GL_n(\OO_k)$ and $t$
diagonal with entries powers of $\varpi$ (\lref{GroundTheory.Automorphic.Cartan.exists_eq_mul_diagPow_mul}). As $\shmiscKone$ is normal in
$\GL_n(\OO_k)$, $\sigma(k_2)w$ is fixed by $\sigma(\shmiscKone)$ (\lref{GlobalTheory.DepthZeroJac.apply_mem_fixed}), and
$\sigma(k_2)w=\chi(\det k_2)\pi(k_2)w$; Step 1 gives $\pi(t)\pi(k_2)w\in V(\shmiscKone)$. Finally $V(\shmiscKone)$ is stable under $\pi(k_1)$, because
$k_1$ normalises $\shmiscKone$ and $\pi(k_1)(\pi(k)z-z)=\pi(k_1kk_1^{-1})\pi(k_1)z-\pi(k_1)z$
(\lref{GlobalTheory.DepthZeroJac.apply_mem_twKer_of_normalizer}, \lref{GlobalTheory.DepthZeroJac.mem_normalizer_KOne}).

\emph{Conclusion.} Since $\pi$ is irreducible and $w\ne0$, the vectors $\pi(g)w$ span $V$ (their span is a non-zero invariant subspace), so
$V(\shmiscKone)=V$ by Step 2; in particular $v\in V(\shmiscKone)$. But for a compact open subgroup $K$ of $\GL_n(k)$ and a smooth representation, a
$K$-fixed vector in $V(K)$ is zero: the projector $e_K$ (averaging over $K$) is the identity on $V^K$ and kills $V(K)$
(\lref{GlobalTheory.IwahoriInj.eq_zero_of_fixed}; $\shmiscKone$ is compact). Hence $v=0$, a contradiction. So $\chi$ is tame.
\end{proof}

\begin{proposition}[Goodness and supports]\label{pf:sh-misc:goodseg}
\begin{enumerate}[label=(\alph*)]
\item Any two supercuspidal supports of an $x\in\shmiscIrr(k)$ are equal.
\item If $x=L(l)$ for a list $l$ of segments, then $x$ is good if and only if $\rho_\Delta$ satisfies \shmiscGo{} for every $\Delta\in l$.
\item A cuspidal $x\in\shmiscIrr(k)$ is good if and only if it satisfies \shmiscGo{}.
\end{enumerate}
\formallabel{UpperTheory.Zero.LLC.goodSeg}
\end{proposition}

\begin{proof}
(a) is the uniqueness of supercuspidal supports \cite[Thm.~2.9]{BernsteinZelevinsky1977}; in the development it is
\lref{GlobalTheory.CompanionTheory.Hypotheses₂.TD.hasSupport_unique}, deduced from the geometric lemma and the finiteness theorem
of Bernstein--Zelevinsky (published results, \lref{PositivityTheory.Exh.geometricLemmaGL}, \lref{UpperTheory.Zero.LLC.Pr.bz}).
Consequently, \emph{if $s$ is any supercuspidal support of $x$, then $x$ is good if and only if every member of $s$ satisfies} \shmiscGo{}:
goodness asks for some support all of whose members satisfy \shmiscGo{}, and all supports equal $s$.

(b) By Zelevinsky's theorem, $L(l)$ has the support $\bigcup_{\Delta\in l}\{\rho_\Delta\nu^j:0\le j<\ell_\Delta\}$ (as a multiset)
(\lref{GlobalTheory.LocGL.hasSupport_msupp_of_isL}, from the published result \lref{UpperTheory.Zero.LLC.Pr.zel},
\cite{Zelevinsky1980}). Condition \shmiscGo{} is invariant under twists by smooth characters, in particular by $\nu^j$
(\ref{pf:llc:F-cl}). Hence every member of the support satisfies \shmiscGo{} if and only if every $\rho_\Delta$ does ($\ell_\Delta\ge1$, so
$\rho_\Delta=\rho_\Delta\nu^0$ is a member).

(c) A cuspidal $x$ has the support $\{x\}$ (\lref{GlobalTheory.LocGL.hasSupport_self}); apply the consequence of (a).
\end{proof}

% ===========================================================================
\subsection{Ray class fields, partial Hecke \texorpdfstring{$L$}{L}-functions and singular moduli}\label{sec:pf:sh-misc:cft}

This subsection proves three classical statements in the form used by the development: the existence of ray class fields
from global Artin reciprocity (used in \cref{sec:pf:baselat2,sec:pf:onto3}), the continuation and non-vanishing on
$\Re s\ge1$ of partial Hecke $L$-functions (used in \cref{sec:pf:bundles,sec:pf:liftinj}), and the integrality of singular
moduli (used in \cref{sec:pf:reckernel,sec:pf:rows}).

Let $K$ be a number field and $\mathfrak m\ne0$ an ideal of $\OO_K$. We write $I^{\mathfrak m}$ for the group of fractional ideals prime to
$\mathfrak m$ and $P_{\mathfrak m}\subseteq I^{\mathfrak m}$ for the subgroup of principal ideals $(\alpha/\beta)$ with $\alpha,\beta\in\OO_K\smallsetminus\{0\}$
prime to $\mathfrak m$, $\alpha\equiv\beta\bmod\mathfrak m$ and $\alpha\beta$ totally positive; $\mathrm{Cl}_{\mathfrak m}=I^{\mathfrak m}/P_{\mathfrak m}$ is the
narrow ray class group modulo $\mathfrak m$. A \emph{global Artin map} for $K$ and an algebraic closure $\Omega$ is a homomorphism
$\mathrm{art}\colon\A_K^\times\to G_K^{\mathrm{ab}}$ which is an open quotient map, whose kernel is the closure of $K^\times(K_\infty^\times)^0$, and
which is compatible with the local inertia groups (the images of the local units at $v$ are the images of the inertia groups of the places
of $\Omega$ above $v$), with Frobenius elements ($\mathrm{art}$ maps a uniformiser at $v$ to the image of $\gamma^{-1}$ for a Frobenius element
$\gamma$ at a place above $v$) and with decomposition groups. \emph{Global Artin reciprocity} is the statement that global Artin maps exist for all
$K$ and $\Omega$ \cite{Tate1967gcft}. The \emph{ray class field statement} is: for every $K$ and every $\mathfrak m\ne0$ there is a finite abelian
extension $L/K$ in which every prime $\mathfrak p\nmid\mathfrak m$ is unramified, whose Galois group is generated by Frobenius elements of
primes $\mathfrak p\nmid\mathfrak m$, and such that an integral ideal $I\ne0$ prime to $\mathfrak m$ lies in $P_{\mathfrak m}$ if and only if
$\prod_{\mathfrak p}\chi(\Frob_{\mathfrak p})^{v_{\mathfrak p}(I)}=1$ for every character $\chi$ of $\Gal(L/K)$.

\begin{lemma}[Two forms of the ray class field statement]\label{pf:sh-misc:rayforms}
The ray class field statement holds if and only if for every $K$ and every $\mathfrak m\ne0$ there are a finite abelian extension $L/K$ and an
isomorphism $\mathrm{Art}\colon\mathrm{Cl}_{\mathfrak m}\to\Gal(L/K)$ such that every prime $\mathfrak p\nmid\mathfrak m$ is unramified in $L$ and
$\mathrm{Art}([\mathfrak p])$ is a Frobenius element at $\mathfrak p$.
\formallabel{GlobalTheory.Published.rayClassFieldSt_iff_artinIso}
\end{lemma}

\begin{proof}
For fixed $K$, $\mathfrak m$ and $L$ the two conditions are compared clause by clause in \lref{GlobalTheory.RayFieldCompare.rayField_iff}:
given the isomorphism, the Frobenius elements of the primes $\mathfrak p\nmid\mathfrak m$ are the images of the classes $[\mathfrak p]$, which
generate $\mathrm{Cl}_{\mathfrak m}$, and the kernel of $I\mapsto\prod_{\mathfrak p}\Frob_{\mathfrak p}^{v_{\mathfrak p}(I)}$ is $P_{\mathfrak m}$, which is
detected by the characters of $\Gal(L/K)$; conversely these properties define a surjection $I^{\mathfrak m}\to\Gal(L/K)$,
$\mathfrak p\mapsto\Frob_{\mathfrak p}$, whose kernel is $P_{\mathfrak m}$ (on integral ideals by assumption, and in general because every class
of the finite group $\mathrm{Cl}_{\mathfrak m}$ contains integral ideals and $J^h\in P_{\mathfrak m}$ for $h=|\mathrm{Cl}_{\mathfrak m}|$). This is the
usual formulation of the existence theorem of class field theory \cite{Neukirch1999}.
\end{proof}

\begin{theorem}[Ray class fields]\label{pf:sh-misc:rayfield}
Global Artin reciprocity implies the ray class field statement.
\formallabel{GlobalTheory.Published.GlobalArtinReciprocity.rayClassFieldSt}
\end{theorem}

\begin{proof}
By \cref{pf:sh-misc:rayforms}, it suffices to prove: \emph{for every $K$ and $\mathfrak m\ne0$ there are a finite abelian extension $L/K$
and an isomorphism $\mathrm{Art}\colon\mathrm{Cl}_{\mathfrak m}\to\Gal(L/K)$ such that every prime $\mathfrak p\nmid\mathfrak m$ is unramified in
$L$ and $\mathrm{Art}([\mathfrak p])$ is a Frobenius element at $\mathfrak p$.} Fix $K$, $\mathfrak m$, an algebraic closure $\Omega$ of $K$ and a global
Artin map $\mathrm{art}$.

Let $U_{\mathfrak m}\subset\A_K^\times$ be the subgroup of ideles $a$ with $a_v\in\OO_v^\times$ for all finite $v$, $a_v\equiv1\bmod\mathfrak m\OO_v$ for
$v\mid\mathfrak m$, and $a_\infty$ in the identity component of $K_\infty^\times$; it is open (\lref{GlobalTheory.RayFromArtin.isOpen_rayIdl}).
Put $R_{\mathfrak m}=K^\times U_{\mathfrak m}$, an open, hence closed, subgroup, and let $\theta\colon\A_K^\times\to\A_K^\times/R_{\mathfrak m}$ be the
quotient map. Since $(K_\infty^\times)^0\subseteq U_{\mathfrak m}$, the kernel of $\mathrm{art}$ is contained in $R_{\mathfrak m}$; as $\mathrm{art}$ is
surjective, there is a unique homomorphism $\rho\colon G_K^{\mathrm{ab}}\to\A_K^\times/R_{\mathfrak m}$ with $\rho\circ\mathrm{art}=\theta$, and $\ker\rho$ is
open because $\mathrm{art}$ is a quotient map and $\mathrm{art}^{-1}(\ker\rho)=R_{\mathfrak m}$. The composite $\varphi\colon G_K\to G_K^{\mathrm{ab}}\to
\A_K^\times/R_{\mathfrak m}$ is surjective with open kernel; by Galois theory there is a finite abelian extension $L\subset\Omega$ of $K$ with
$G_L=\ker\varphi$ (\lref{GlobalTheory.RayFromArtin.exists_field}), and $\varphi$ induces an isomorphism
$e\colon\A_K^\times/R_{\mathfrak m}\to\Gal(L/K)$ with $e(\varphi(\sigma))=\sigma|_L$.

For $J\in I^{\mathfrak m}$ let $u(J)$ be the idele with component $\varpi_{\mathfrak p}^{v_{\mathfrak p}(J)}$ at every $\mathfrak p$ (fixed uniformisers) and
$1$ elsewhere, and put $\psi(J)=e(\theta(u(J)))^{-1}$; $\psi\colon I^{\mathfrak m}\to\Gal(L/K)$ is a homomorphism.

\emph{Unramified primes.} Let $\mathfrak p\nmid\mathfrak m$, $A$ a place of $\Omega$ above $\mathfrak p$ and $\sigma\in I_A$. By the inertia compatibility,
the image of $\sigma$ in $G_K^{\mathrm{ab}}$ is $\mathrm{art}(x)$ for an idele $x$ concentrated at $\mathfrak p$ with unit component; such an $x$ lies in
$U_{\mathfrak m}$ (its component at $\mathfrak p$ is a unit, its components at the
places dividing $\mathfrak m$ are $1$, and its archimedean component is $1$). Hence $\varphi(\sigma)=\theta(x)=1$ and $\sigma|_L=1$. Since every element
of the inertia group of a prime $\mathfrak P\mid\mathfrak p$ of $L$ lifts to an element of $I_A$ for a place $A$ of $\Omega$ above $\mathfrak P$ (an element
$\tau$ with $\tau x-x\in\mathfrak P$ for $x\in\OO_L$ acts trivially on the residue field of $A\cap L$, as one sees by writing elements of
the valuation ring as quotients $t/s$ with $s\notin\mathfrak P$), the inertia groups of $L/K$ at $\mathfrak p$ are trivial, i.e.\ $\mathfrak p$ is
unramified in $L$.

\emph{Frobenius elements.} For $\mathfrak p\nmid\mathfrak m$ and a place $A$ above $\mathfrak p$, the Frobenius compatibility gives a Frobenius element
$\gamma$ at $A$ with $\mathrm{art}(u(\mathfrak p))=\gamma^{-1}$ in $G_K^{\mathrm{ab}}$. Then $\psi(\mathfrak p)=e(\varphi(\gamma^{-1}))^{-1}=\gamma|_L$, a Frobenius
element of $L/K$ at $\mathfrak p$ (\lref{GlobalTheory.RayFromArtin.isFrob_restrict}).

\emph{Kernel.} $J\in\ker\psi$ if and only if $u(J)\in R_{\mathfrak m}$, and this holds if and only if $J\in P_{\mathfrak m}$
(\lref{GlobalTheory.RayFromArtin.unifIdl_mem_rayGrp_iff}, the usual comparison between ideles and ideals).

\emph{Surjectivity.} The Frobenius elements of the primes $\mathfrak p\nmid\mathfrak m$ generate $\Gal(L/K)$
(\lref{GlobalTheory.RayMaximal.closure_isFrob_eq_top}, a consequence of the Chebotarev density theorem), and each of them is
$\psi(\mathfrak p)$, since in an abelian extension the Frobenius element at an unramified prime is unique.

Hence $\psi$ induces an isomorphism $\mathrm{Art}\colon\mathrm{Cl}_{\mathfrak m}=I^{\mathfrak m}/P_{\mathfrak m}\to\Gal(L/K)$ with the required properties.
\end{proof}

\begin{remark}\label{pf:sh-misc:rayrem}
The Artin isomorphism constructed in the proof is normalised so that classes of primes go to (arithmetic) Frobenius elements; the
formal definition of a global Artin map uses the opposite normalisation for uniformisers, which is compensated by the inverse in the
definition of $\psi$. The ray class group is the narrow one (positivity at all real places).
\end{remark}

For a number field $K$, a finite set $T$ of finite places and a Hecke character $\chi$ (a character of $\A_K^\times/K^\times$, with value
$\chi(\mathfrak p)$ at unramified $\mathfrak p$), the partial $L$-function is $L^T(s,\chi)=\prod_{\mathfrak p\notin T}(1-\chi(\mathfrak p)\Nm\mathfrak p^{-s})^{-1}$
for $\Re s>1$, the product running over the primes at which $\chi$ is unramified; $\zeta_K^T=L^T(\cdot,1)$.

\begin{theorem}[Partial Hecke $L$-functions]\label{pf:sh-misc:heckeL}
Let $K$ be a number field and $T$ a finite set of finite places of $K$.
\begin{enumerate}[label=(\roman*)]
\item For every unitary Hecke character $\chi$ of $K$ which is not of the form $|\cdot|^{it}$, $t\in\R$, there is an entire function $\Lambda$
with $\Lambda(s)=L^T(s,\chi)$ for $\Re s>1$ and $\Lambda(s)\ne0$ for $\Re s\ge1$.
\item There is an entire function $Z$ with $Z(s)=(s-1)\zeta_K^T(s)$ for $\Re s>1$ and $Z(s)\ne0$ for $\Re s\ge1$.
\end{enumerate}
\formallabel{GlobalTheory.Published.HeckeLPartialContinuation_holds}
\end{theorem}

This is Hecke's theorem \cite{Hecke1917,Hecke1920} together with the non-vanishing of Hecke $L$-functions on the line $\Re s=1$
(see e.g.\ \cite[Ch.~VII]{Neukirch1999}); the removal of finitely many Euler factors does not affect these properties. The formal
development proves it as follows.

\begin{proof}
For a finite set $T$ and $\Re s>0$ put $r_T(\chi,s)=\prod_{\mathfrak p\in T}(1-c_\chi(\mathfrak p)\Nm\mathfrak p^{-s})$, where $c_\chi(\mathfrak p)=\chi(\mathfrak p)$
if $\chi$ is unramified at $\mathfrak p$ and $0$ otherwise; it is entire in $s$, and non-zero for $\Re s>0$ because
$|c_\chi(\mathfrak p)\Nm\mathfrak p^{-s}|<1$ for unitary $\chi$. For $\Re s>1$, $L^T(s,\chi)=L(s,\chi)\,r_T(\chi,s)$, where $L(s,\chi)$ is the full Euler product
(\lref{GlobalTheory.HeckePartial.heckeL_eq_mul_removed}).

(i) There is an entire $\Lambda_0$ equal to the Euler product $L(s,\chi)$ for $\Re s>1$
(\lref{GlobalTheory.HeckeOne.entire_of_not_normPow}), and $\Lambda_0(s)\ne0$ for $\Re s\ge1$
(\lref{GlobalTheory.HeckePartial.IsHeckeL.ne_zero_of_one_le}). Take $\Lambda=\Lambda_0\,r_T(\chi,\cdot)$.

(ii) There is an entire $Z_0$ with $Z_0(1)\ne0$ and $Z_0(s)=(s-1)\exp\bigl(-\sum_{\mathfrak p}\log(1-\Nm\mathfrak p^{-s})\bigr)=(s-1)\zeta_K(s)$ for $\Re
s>1$ (\lref{GlobalTheory.dedekindZeta_eulerProduct_mul_sub_one_entire_ne_zero}). Take $Z=Z_0\,r_T(1,\cdot)$; then
$Z(s)=(s-1)\zeta_K^T(s)$ for $\Re s>1$. For $\Re s>1$, $Z_0(s)$ is a product of $s-1\ne0$ and an exponential; $Z_0(1)\ne0$; and
$Z_0(1+it)\ne0$ for real $t\ne0$ by the non-vanishing of Euler products of this type on the line $\Re s=1$
(\lref{GlobalTheory.zetaEulerProduct_ne_zero_one_add_mul_I}, applied to the norms $\Nm\mathfrak p\ge2$, for which
$\sum_{\mathfrak p}\Nm\mathfrak p^{-\sigma}$ converges for $\sigma>1$). Since $r_T(1,s)\ne0$ for $\Re s>0$, $Z(s)\ne0$ for $\Re s\ge1$.
\end{proof}

\begin{remark}\label{pf:sh-misc:heckeLrem}
The development labels this statement as a published result; it is nevertheless proved there, from the two cited continuation statements
and the non-vanishing statements on $\Re s=1$.
\end{remark}

For a lattice $L\subset\C$ with invariants $g_2(L),g_3(L)$ put $j(L)=1728g_2^3/(g_2^3-27g_3^2)$, and for a positive integer $d$ let
$\omega_d=(1+i\sqrt d)/2$. We say that $L$ \emph{has complex multiplication by} $d$ if $\omega_dL\subseteq L$.

\begin{theorem}[Integrality of singular moduli]\label{pf:sh-misc:cm1}
Let $d$ be a squarefree positive integer with $d\equiv3\bmod4$ and $L\subset\C$ a lattice with complex multiplication by $d$. Then $j(L)$ is
integral over $\Z$.
\formallabel{ReciprocityTheory.ChartRestriction.L1.cm1}
\end{theorem}

This is the classical integrality of singular moduli \cite[Ch.~II, Thm.~6.1]{Silverman1994}, \cite[\S11]{Cox1989}. The formal proof
follows the classical argument via the modular equation:

\begin{proof}
Write $L=c(\Z\tau+\Z)$ with $c\in\C^\times$ and $\tau$ in the upper half plane (\lref{ReciprocityTheory.ChartRestriction.L1.latTau});
then $j(L)=j(\tau)$ (\lref{ReciprocityTheory.ChartRestriction.L1.jInv_eq_jT}). Since $\alpha L\subseteq L$ for $\alpha=x+y\omega_d$
($x,y\in\Z$), there are integers $A,B,C,D$ with $\alpha\tau=A\tau+B$, $\alpha=C\tau+D$ and
$4(AD-BC)=4x^2+4xy+y^2(1+d)$ (\lref{ReciprocityTheory.ChartRestriction.L1.cmPoint}: the matrix of multiplication by $\alpha$ in the basis
$(\tau,1)$ has determinant $\Nm\alpha$). Take $(x,y)=(-1,2)$, i.e.\ $\alpha=-1+2\omega_d=i\sqrt d$: then $AD-BC=d$ and $A\tau+B=\tau(C\tau+D)$, so
$\tau$ is a fixed point of an integral matrix of determinant $d$. Let
\[
G_d(\tau)=\prod_{ad'=d,\ 0\le b<d'}\Bigl(\mathcal A(\tau)\Delta\bigl(\tfrac{a\tau+b}{d'}\bigr)-\mathcal A\bigl(\tfrac{a\tau+b}{d'}\bigr)\Delta(\tau)\Bigr),
\]
where $j=\mathcal A/\Delta$ with $\mathcal A=1728g_2^3$ and $\Delta$ the discriminant (as functions of $\tau$). A fixed point of an integral matrix of
determinant $d$ makes one factor vanish, so $G_d(\tau)=0$ (\lref{ReciprocityTheory.ChartRestriction.L1.Gm_vanish}). Since $d\equiv3\bmod4$, $d$
is not a square, and $G_d(\tau)=0$ implies that $j(\tau)$ is integral over $\Z$
(\lref{ReciprocityTheory.ChartRestriction.L1.jT_isIntegral}: up to a non-zero factor, $G_d(\tau)$ is the value at $X=j(\tau)$ of the
polynomial $\Phi_d(X,X)$, which for non-square $d$ has integer coefficients and leading coefficient $\pm1$, \cite[\S11]{Cox1989}).
\end{proof}

% ===========================================================================
\subsection{Intersection numbers of multiprojective varieties}\label{sec:pf:sh-misc:intnum}

This subsection proves that the intersection numbers of the development form a field intersection family, and that field intersection
families exist and are unique on projective varieties; these facts are used in \cref{sec:pf:deepstr,sec:pf:towergeo}. We use the
multigraded polynomial ring $F[x]=F[x_{k,i}:k\in K,i\in\iota_k]$ over a field $F$ (with a finite set $K$ of blocks), projective varieties,
charts, rational sections, divisors $\mathrm{div}_V(f/g)$ and field intersection families as in \cref{def:pf:deepstr:fieldint}. In axiom (i)
of a field intersection family the degree of a point $V$ (a projective variety of dimension $0$) is computed in a chart $i=(i_k)_k$
(variables $x_{k,i_k}\notin V$): $I(V;\varnothing)=\deg_i(V)$, the point degree of $V$ in the chart $i$.

\begin{definition}[Intersection numbers]\label{pf:sh-misc:intnumdef}
Let $P\subseteq F[x]$ be an ideal. For $m\in\Z^K$ let $h_P(m)=\dim_F(F[x]/P)_m$, the dimension of the image of the polynomials of
multidegree $m$, if all $m_k\ge0$, and $h_P(m)=0$ otherwise. For $b\in\Z^K$ and $h\colon\Z^K\to\Z$ put $(\Delta_bh)(m)=h(m)-h(m-b)$.
For $b_1,\dots,b_n\in\Z^K$ the \emph{intersection number} $\mathrm{int}_F(P;b_1,\dots,b_n)$ is the eventual value of
$\Delta_{b_1}\cdots\Delta_{b_n}h_P(m)$ as $m\to\infty$ (in the product order of $\Z^K$): the constant $c$ with
$\Delta_{b_1}\cdots\Delta_{b_n}h_P(m)=c$ for all large $m$, if it exists, and $0$ otherwise.
\formal{GroundTheory.Heights₂.Heights.NumbersDVR.intNum}
\end{definition}

These are the mixed multiplicities of $F[x]/P$, read off from its multigraded Hilbert function.
We use two results proved elsewhere:
\begin{enumerate}[label=(\alph*)]
\item (\lref{GroundTheory.Heights₂.Heights.FieldNumbers.IntNumBase.full}) If $P$ is a multiprojective prime of relative dimension $0$,
then the degree-zero part of the homogeneous fraction field of $F[x]/P$ is finite-dimensional over $F$ and $\mathrm{int}_F(P;\varnothing)$
is its dimension; consequently (\lref{GroundTheory.Heights₂.Heights.NumbersDVR.IntNumPoint.proof}) $\mathrm{int}_F(V;\varnothing)=
\deg_i(V)$ for every projective variety $V$ of dimension $0$ and every chart $i$ of $V$.
\item (\lref{GroundTheory.Heights₂.Heights.FieldNumbers.CutRational.full}) If $P$ is a multiprojective prime of relative dimension $n+1$
and $f,g\notin P$ are multihomogeneous of multidegrees $b_f,b_g\in\N^K$, then for all $b_1,\dots,b_n$,
$\mathrm{int}_F(P;b_1,\dots,b_n,b_f-b_g)=\sum_W\mathrm{ord}_W(f/g)\,\mathrm{int}_F(W;b_1,\dots,b_n)$; consequently
(\lref{GroundTheory.Heights₂.Heights.NumbersDVR.IntNumRatSection.proof}) the same holds for every projective variety $V$ of
dimension $n+1$ and every rational section $f/g$ on $V$.
\end{enumerate}

\begin{theorem}[Intersection numbers form a field intersection family]\label{pf:sh-misc:intnumfield}
For every field $F$, every finite set $K$ and all finite sets $\iota_k$, the intersection numbers $(V;L_1,\dots,L_n)\mapsto\mathrm{int}_F(V;L_1,
\dots,L_n)$ form a field intersection family over $F$.
\formallabel{GroundTheory.Heights₂.Heights.NumbersDVR.IntNumIsFieldIntersection.full}
\end{theorem}

\begin{proof}
Axiom (i) is (a), and axiom (ii) is (b): for a projective variety $V$ of dimension $n+1$ and a rational section $f/g$ of multidegree
$A=b_f-b_g$, $\mathrm{int}_F(V;L,A)=\sum_W\mathrm{ord}_W(f/g)\,\mathrm{int}_F(W;L)$. This deduction also identifies the notions
of multihomogeneous elements and ideals and of divisors used in (a), (b) with those of \cref{def:pf:deepstr:fieldint}
(\lref{GroundTheory.Heights₂.Heights.NumbersDVR.VocabMultiHomog.proof},
\lref{GroundTheory.Heights₂.Heights.NumbersDVR.VocabMultiHomogIdeal.proof},
\lref{GroundTheory.Heights₂.Heights.NumbersDVR.VocabDivCycle.proof}); and a projective variety is a multiprojective prime with
the stated Krull dimension, which is its relative dimension plus $|K|$.
\end{proof}

\begin{theorem}[Existence and uniqueness]\label{pf:sh-misc:existsunique}
For every field $F$, every finite set $K$ and all finite sets $\iota_k$ there is a field intersection family $I$ over $F$ such that every
field intersection family $J$ over $F$ satisfies $J(V;L)=I(V;L)$ for every projective variety $V$ of dimension $n$ and all
$L\in(\Z^K)^n$.
\formallabel{GroundTheory.Heights₂.Heights.Targets.ExistsUnique.full}
\end{theorem}

\begin{proof}
Take $I=\mathrm{int}_F$, a field intersection family by \cref{pf:sh-misc:intnumfield}. Let $J$ be a field intersection family. Then $J$ is an
intersection theory in the sense of \lref{GroundTheory.Heights₂.Heights.FieldNumbers.IntNumUnique.full}: by axiom (i) and (a) it
takes the value $\dim_F$ of the degree-zero homogeneous fraction field at multiprojective primes of relative dimension $0$, and by axiom (ii)
it satisfies the divisor formula of (b) (\lref{GroundTheory.Heights₂.Heights.NumbersDVR.Numbers.isIntersectionTheory}). By the
uniqueness statement \lref{GroundTheory.Heights₂.Heights.FieldNumbers.IntNumUnique.full} (an intersection theory agrees with
$\mathrm{int}_F$ at every multiprojective prime of relative dimension $n$ and all $L\in(\Z^K)^n$), $J(V;L)=\mathrm{int}_F(V;L)$ for every projective variety $V$ of dimension $n$
(\lref{GroundTheory.Heights₂.Heights.NumbersDVR.Numbers.agree}).
\end{proof}

% ===========================================================================
\subsection{The second Lie derivative formula on the Fock model}\label{sec:pf:sh-misc:lie2}

This subsection proves the infinitesimal form of the action of the second unitary group of a Weil representation on the
Fock model, used in \cref{sec:pf:thmK} (hypothesis \ref{pf:thmK:H:lie2}) and in \cref{sec:pf:tlift} (\cref{pf:tlift:fock}).
The formal statement is an implication: the existence of local metaplectic lifts on Schwartz space and the existence of the
Heisenberg derivatives imply the formula. Both hypotheses are theorems of the development
(\lref{PositivityTheory.ThKR.localLift_holds} for the first). The proof determines the scalar part of the derivative by a
doubling argument; the identification of the derivative with a quadratic operator up to a scalar, and the Fock-model
computation, are two results of the formal development which we state precisely and cite (\cref{pf:sh-misc:main2,pf:sh-misc:fockdata}).


\paragraph{Setting.}
We use the notation of \cref{sec:pf:tlift} (Schwartz--Bruhat functions, Darboux coordinates $e=(e_1,e_2)$, Heisenberg operators
$\rho$, admissible linear forms, pure tensors $f\otimes g$ with $f\in\shmiscS(F_\infty^\kappa)$ and $g$ a Schwartz--Bruhat function on
$\A_f^\kappa$). Let $F\subset E$ be a CM extension with conjugation $\sigma$, $\psi$ the standard additive character of $\A_F$,
$\tau=(\tau_w)_w$ a choice of embeddings $E\to\C$ above the infinite places $w$ of $F$, and $J_A\in M_\alpha(\A_E)$, $J_B\in M_\beta(\A_E)$.
The space $W=M_{\alpha\times\beta}(\A_E)$ carries the $\A_F$-valued alternating form $\langle\cdot,\cdot\rangle_\A$ built from
$J_A,J_B$ (as in \cref{sec:pf:tlift}), and $d$ denotes Darboux coordinates for it: an $\A_F$-linear isomorphism
$e=(e_1,e_2)\colon W\to\A_F^\kappa\times\A_F^\kappa$ with $\langle u,v\rangle_\A=e_1(u)\cdot e_2(v)-e_1(v)\cdot e_2(u)$. We write
$\langle u,v\rangle=-\Tr_{F_\infty/\R}\bigl((\langle u,v\rangle_\A)_\infty\bigr)$. An element of $W$ (or of $\A_E$, or a
matrix) is \emph{archimedean} if it is supported at infinity. The \emph{doubled} data are $J_B^{(2)}=\mathrm{diag}(J_B,-J_B)$, the space
$W^{(2)}=M_{\alpha\times2\beta}(\A_E)$ of pairs $(u\,|\,u')$ of blocks of columns, with the form built from $J_A,J_B^{(2)}$, and doubled
Darboux coordinates $d^{(2)}$, given by $e^{(2)}(u\,|\,u')=\bigl((e_1(u),e_1(u')),(e_2(u),-e_2(u'))\bigr)$. For $H\in M_\beta(\C)$, a
\emph{frame} of $H$ of signature $(R,S)$ is a matrix $C\in M_{\beta\times(R+S)}(\C)$ which has a right inverse and satisfies
${}^tC\,H\,\bar C=J_{R,S}=\mathrm{diag}(1_R,-1_S)$. For a family $(T(t))_{t\in\R}$ of operators on functions on $\A^\kappa$
and $f,f_1\in\shmiscS(F_\infty^\kappa)$ we say that \emph{$T(t)f$ has derivative $f_1$ at $0$} if for every Schwartz--Bruhat function
$g$ on $\A_f^\kappa$ and every admissible linear form $\ell$ the function $t\mapsto\ell(T(t)(f\otimes g))$ is differentiable at $0$ with
derivative $\ell(f_1\otimes g)$. A \emph{one-parameter subgroup} of $\Uni(J_B)(\A)$ is a continuous homomorphism $\gamma\colon\R\to
\Uni(J_B)(\A)$ such that $\gamma(t)-1$ is archimedean for all $t$.

\begin{definition}[The three statements]\label{pf:sh-misc:liedefs}
\begin{enumerate}[label=(\alph*)]
\item $\sfname{LocalLift}$: for every finite set $l$ and every $X\in\mathfrak{sp}_{2l}(\R)$ there are $\eps>0$, families $V(t),W(t)$ of
continuous operators on $\shmiscS(\R^l)$ ($|t|<\eps$) and $c\in\C$ such that $V(0)=1$, $V(t)W(t)=W(t)V(t)=1$, $V(t)$ implements
$\exp(tX)$ in the Schr\"odinger model, $t\mapsto\lambda(V(t)f)$ is differentiable for $|t|<\eps$ and every continuous linear form
$\lambda$, and its derivative at $0$ is $\lambda(q(X)f+cf)$, where $q(X)$ is the quadratic operator attached to $X$ by the
oscillator representation.
\item $\sfname{DerEx}$: for all data as above there are continuous operators $D_u$ on $\shmiscS(F_\infty^\kappa)$, $u\in W$ archimedean,
such that $\rho(tu)f$ has derivative $D_uf$ at $0$ and
\[
(D_uf)(x)=\partial f(x)\bigl[e_1(u)_\infty\bigr]-2\pi i\,\Tr_{F_\infty/\R}\bigl(x\cdot e_2(u)_\infty\bigr)f(x).
\]
\item $\sfname{LieSecond}$: for all data as above, all $a,b\in\C$ with $ab=i$ and all families of frames $A_w$ of $a\tau_w(J_A)$ (of
signature $(P,Q)$) and $C_w$ of $b\tau_w(J_B)$ (of signature $(R,S)$), every Fock map $\Phi\colon\mathcal T(P,Q,R,S)\to\shmiscS(F_\infty^\kappa)$
for $(\psi,\tau,d,A,C)$, every unitary Hecke character $\chi$ of $E$ with \emph{infinity type} $m=(m_w)\in\Z^{\{w\}}$, i.e.\
$\chi(z)=\prod_w(\tau_w(z)/|\tau_w(z)|)^{m_w}$ for every archimedean idele $z$ with $z-1$ archimedean, and every representation $\omega$ of
$\Uni(J_B)(\A)$ which is a Weil representation for $J_B$ split by $\chi$ (conditions (W1) and (W2) of \cref{pf:tlift:weil} for the
group $\Uni(J_B)$, with doubled Darboux coordinates $d^{(2)}$ and doubled representation $\omega^{(2)}$ of $\Uni(J_B^{(2)})(\A)$,
$J_B^{(2)}=\mathrm{diag}(J_B,-J_B)$): for every one-parameter subgroup $\gamma$ of $\Uni(J_B)(\A)$ with generators $Y_w$ in the frames,
i.e.\ $\tau_w(\gamma(t))C_w=C_w\exp(tY_w)$, and every $T\in\mathcal T(P,Q,R,S)$, $\omega(\gamma(t))\Phi(T)$ has derivative
\[
\Phi\Bigl(\sum_w\bigl(\shmiscR(Y_w)+\tfrac{m_w}2\tr Y_w\bigr)_wT\Bigr)
\]
at $0$, where $\shmiscR$ is the second Fock action and $(\cdot)_w$ means acting on the $w$-th factor.
\end{enumerate}
\formal{PositivityTheory.ThKR.LocalLift}\formal{GlobalTheory.RealLift.DerEx}\formal{GlobalTheory.RealLift.LieSecond}
\formal{Formalism.BasicNotions.TH.fockAct2}
\end{definition}

The second Fock action and the Fock maps are those of the formal development; in the definite case they give the formulas of
\cref{pf:tlift:fock}.

\begin{theorem}[The second Lie derivative formula]\label{pf:sh-misc:lie2}
$\sfname{LocalLift}$ and $\sfname{DerEx}$ imply $\sfname{LieSecond}$.
\formallabel{PositivityTheory.ThKR.lieSecond2}
\end{theorem}

We use the following results of the formal development. For a one-parameter subgroup $\gamma$ with $\tau_w(\gamma(t))=\exp(tZ_w)$, a
\emph{decomposition} of $Z=(Z_w)$ is a finite family $(u_j,v_j)_j$ of archimedean elements of $W$ with
\begin{equation}\label{eq:pf:sh-misc:decomp}
\tau_w(x)\,{}^tZ_w=\sum_j\langle v_j,x\rangle\,\tau_w(u_j)\qquad\text{for every archimedean }x\in W\text{ and every }w .
\end{equation}

\begin{lemma}[Derivatives of Weil representations up to a scalar]\label{pf:sh-misc:main2}
Assume $\sfname{LocalLift}$. Let the data $F,E,\psi,\tau,J_A,J_B,d$ be as above, let $K_w$ be invertible hermitian matrices such that
$g^*K_wg=K_w$ if and only if ${}^tg\,\tau_w(J_B)\,\bar g=\tau_w(J_B)$, let $\omega$ be a representation of $\Uni(J_B)(\A)$ satisfying
(W1), and let $(D_u)$ be as in $\sfname{DerEx}$. Then there are $c_w\in\C$ such that for every one-parameter subgroup $\gamma$ with
$\tau_w(\gamma(t))=\exp(tZ_w)$, every decomposition $(u_j,v_j)$ of $Z$ and every $f\in\shmiscS(F_\infty^\kappa)$, $\omega(\gamma(t))f$ has
derivative
\[
(8\pi i)^{-1}\sum_j\bigl(D_{u_j}D_{v_j}f+D_{v_j}D_{u_j}f\bigr)+\Bigl(\sum_wc_w\tr Z_w\Bigr)f
\]
at $0$. \textup{(}\lref{PositivityTheory.ThKR.main2}.\textup{)}
\end{lemma}

\begin{lemma}[Fock data]\label{pf:sh-misc:fockdata}
Let $A,C$ be frames, $\Phi$ a Fock map, $(D_u)$ as in $\sfname{DerEx}$, and $\gamma$ a one-parameter subgroup with generators $Y_w$ in the frame
$C$. Then there is $Z=(Z_w)$ with $\tr Z_w=\tr Y_w$, $\tau_w(\gamma(t))=\exp(tZ_w)$, $\tau_w(\det\gamma(t))=\exp(t\tr Y_w)$ and
$\Re\tr Y_w=0$, and there is a decomposition $(u_j,v_j)$ of $Z$ such that for every $T$
\[
(8\pi i)^{-1}\sum_j\bigl(D_{u_j}D_{v_j}+D_{v_j}D_{u_j}\bigr)\Phi(T)=\Phi\Bigl(\sum_w\shmiscR(Y_w)_wT\Bigr).
\]
\textup{(}\lref{PositivityTheory.RL.Up.fock_data}.\textup{)}
\end{lemma}

\begin{lemma}[Frames]\label{pf:sh-misc:frames}
\begin{enumerate}[label=(\alph*)]
\item If $C$ is a frame of $H\in M_\beta(\C)$, then $\bar H$ is hermitian and invertible.
\item Let $b\in\C$ and $M\in M_\beta(\C)$ be such that $K=\overline{bM}$ is invertible. Then for every $g\in M_\beta(\C)$: $g^*Kg=K$ if and
only if ${}^tg\,M\,\bar g=M$.
\end{enumerate}
\formallabel{PositivityTheory.RL.Up.frame_K}
\formallabel{PositivityTheory.RL.Up.iff_K}
\end{lemma}

\begin{proof}
(a) Let ${}^tC\,H\,\bar C=J_{R,S}$ and $CC'=1$. As $J_{R,S}^2=1$, ${}^tC$ has the right inverse $H\bar CJ_{R,S}$, so $C$ has rank $R+S$; as
$CC'=1$, $C$ has rank $|\beta|$. Hence $C$ is square and invertible, with inverse $C'$. Conjugating the frame equation gives
$C^*\bar HC=J_{R,S}$, so $\bar H=C'^*J_{R,S}C'$, which is hermitian and invertible (with inverse $CJ_{R,S}C^*$).
(b) If $\beta=\varnothing$ there is nothing to prove; otherwise $b\ne0$, as $K$ is invertible. Since $\overline{{}^tg\,M\bar g}=g^*\bar Mg$, we have
$g^*Kg=\overline{b\,{}^tg\,M\bar g}$, and the assertion follows because complex conjugation is injective and $b\ne0$.
\end{proof}

\begin{lemma}[Doubling]\label{pf:sh-misc:dbl}
\begin{enumerate}[label=(C\arabic*)]
\item For one-parameter subgroups $\gamma,\gamma'$ of $\Uni(J_B)(\A)$, $t\mapsto\mathrm{diag}(\gamma(t),\gamma'(t))$ is a one-parameter subgroup of
$\Uni(J_B^{(2)})(\A)$.
\item $\langle(u\,|\,u'),(v\,|\,v')\rangle^{(2)}=\langle u,v\rangle-\langle u',v'\rangle$ for $u,u',v,v'\in W$, where $\langle\cdot,\cdot\rangle^{(2)}$
is the real form of the doubled space.
\item Let $(D_u)$ and $(D^{(2)}_u)$ satisfy the formula of $\sfname{DerEx}$ for $d$ and $d^{(2)}$, let $u\in W$ be archimedean, and let
$B(x_1,x_2)=f(x_1)f'(x_2)$ with $f,f'\in\shmiscS(F_\infty^\kappa)$ and $B\in\shmiscS(F_\infty^{2\kappa})$. Then
$D^{(2)}_{(u|0)}B(x_1,x_2)=(D_uf)(x_1)\,f'(x_2)$.
\item There is a real-valued $f\in\shmiscS(F_\infty^\kappa)$ with compact support and $f(0)=1$ such that $(x_1,x_2)\mapsto f(x_1)f(x_2)$ lies
in $\shmiscS(F_\infty^{2\kappa})$.
\end{enumerate}
\formallabel{PositivityTheory.RL.Up.exists_dbl}
\formallabel{PositivityTheory.RL.Up.sB_dbl}
\formallabel{PositivityTheory.RL.Up.box_D}
\formallabel{PositivityTheory.RL.Up.exists_box}
\end{lemma}

\begin{proof}
(C1) $\delta(t)=\mathrm{diag}(\gamma(t),\gamma'(t))$ is a continuous homomorphism, and ${}^t\delta(t)\,J_B^{(2)}\,\sigma(\delta(t))
=\mathrm{diag}\bigl({}^t\gamma(t)J_B\sigma(\gamma(t)),-{}^t\gamma'(t)J_B\sigma(\gamma'(t))\bigr)=J_B^{(2)}$; and $\delta(t)-1=\mathrm{diag}(\gamma(t)-1,
\gamma'(t)-1)$ is archimedean. (C2) By the Darboux property of $d^{(2)}$ and the formula for $e^{(2)}$,
$\langle(u|u'),(v|v')\rangle^{(2)}_\A=e_1(u)\cdot e_2(v)-e_1(u')\cdot e_2(v')-e_1(v)\cdot e_2(u)+e_1(v')\cdot e_2(u')=\langle u,v\rangle_\A-\langle
u',v'\rangle_\A$; apply $-\Tr_{F_\infty/\R}$ at infinity. (C3) The element $(u|0)$ is archimedean and $e^{(2)}(u|0)=((e_1(u),0),(e_2(u),0))$.
By the product rule, the derivative of $B$ at $(x_1,x_2)$ in the direction $(e_1(u)_\infty,0)$ is $\partial f(x_1)[e_1(u)_\infty]\,f'(x_2)$,
and $(x_1,x_2)\cdot(e_2(u)_\infty,0)=x_1\cdot e_2(u)_\infty$. Inserting this into the formula of $\sfname{DerEx}$ for $D^{(2)}_{(u|0)}$ gives
$\bigl(\partial f(x_1)[e_1(u)_\infty]-2\pi i\Tr_{F_\infty/\R}(x_1\cdot e_2(u)_\infty)f(x_1)\bigr)f'(x_2)=(D_uf)(x_1)f'(x_2)$. (C4) Take a smooth
bump function $f$ on the real vector space $F_\infty^\kappa$, with values in $[0,1]$, equal to $1$ on the unit ball and supported in the
ball of radius $2$; $f$ and $f\otimes f$ are smooth with compact support, hence Schwartz functions.
\end{proof}

We also use: (C5) there are a Schwartz--Bruhat function $g$ on $\A_f^\kappa$ and $y\in\A_f^\kappa$ with $g(y)=1$
(\lref{PositivityTheory.RL.SeminormBoundRestriction.exists_SB}); (C6) there is an admissible linear form $\ell$ on Schwartz--Bruhat
functions on $\A^{\kappa}\times\A^{\kappa}$ which is invariant under the Heisenberg operators $\rho^{(2)}$ of the diagonal elements
$(z\,|\,z)$, $z\in W$, together with $B\in\shmiscS(F_\infty^{2\kappa})$ and a Schwartz--Bruhat $G$ with $\ell(B\otimes G)\ne0$
(\lref{PositivityTheory.RL.Up.exists_form}); (C7) the operators of $\sfname{DerEx}$ satisfy the canonical commutation
relation $D_uD_vf=D_vD_uf+2\pi i\langle u,v\rangle f$ for archimedean $u,v$ (\lref{PositivityTheory.RL.Up.IsD.comm}), and likewise
for the doubled space.

\begin{proof}[Proof of \cref{pf:sh-misc:lie2}]
Let the data of $\sfname{LieSecond}$ be given; let $d^{(2)},\omega^{(2)}$ be the doubling data of $\omega$ (so $\omega^{(2)}$ satisfies the
Siegel rule with $\chi$, and $\omega^{(2)}(\mathrm{diag}(h,1))(\Phi_1\otimes\Phi_2)=\omega(h)\Phi_1\otimes\Phi_2$), $(D_u)$ and $(D^{(2)}_u)$ the
operators of $\sfname{DerEx}$ for $d$ and $d^{(2)}$. Put $K_w=\overline{b\tau_w(J_B)}$. Since $C_w$ is a frame of $b\tau_w(J_B)$,
\cref{pf:sh-misc:frames} shows that $K_w$ is hermitian and invertible and that $g^*K_wg=K_w$ if and only if ${}^tg\,\tau_w(J_B)\,\bar
g=\tau_w(J_B)$; so the $K_w$ are as required in \cref{pf:sh-misc:main2} for $J_B$. Likewise $\overline{b\tau_w(J_B^{(2)})}=\mathrm{diag}(K_w,-K_w)$
is hermitian and invertible, and \cref{pf:sh-misc:frames}(b) applies to it, so these matrices are as required for $J_B^{(2)}$. Let $(c_w)$ and $(c'_w)$ be the constants of \cref{pf:sh-misc:main2} for $\omega$ and for $\omega^{(2)}$. Let $\gamma$, $Y$, $T$ be
given, and let $Z$ and the decomposition $(u_j,v_j)$ be as in \cref{pf:sh-misc:fockdata}. By \cref{pf:sh-misc:main2,pf:sh-misc:fockdata},
$\omega(\gamma(t))\Phi(T)$ has derivative
\begin{equation}\label{eq:pf:sh-misc:lieder}
\Phi\Bigl(\sum_w\shmiscR(Y_w)_wT\Bigr)+\Bigl(\sum_wc_w\tr Y_w\Bigr)\Phi(T)
\end{equation}
at $0$. It therefore suffices to show $\sum_wc_w\tr Y_w=\sum_w\frac{m_w}2\tr Y_w$, which we do in two steps.

\emph{Step 1: $\sum_wc_w\tr Y_w=\sum_wc'_w\tr Y_w$.} By (C1), $\delta(t)=\mathrm{diag}(\gamma(t),1)$ is a one-parameter subgroup of
$\Uni(J_B^{(2)})(\A)$, with $\tau_w(\delta(t))=\exp(t\,\mathrm{diag}(Z_w,0))$. The family $\bigl((u_j|0),(v_j|0)\bigr)_j$ is a decomposition of
$\mathrm{diag}(Z,0)$: for $x=(x_1|x_2)$, (C2) gives $\langle(v_j|0),x\rangle^{(2)}=\langle v_j,x_1\rangle$, and $(x_1|x_2)\,{}^t\mathrm{diag}(Z,0)=
(x_1{}^tZ\,|\,0)=\sum_j\langle v_j,x_1\rangle(u_j|0)$ by \eqref{eq:pf:sh-misc:decomp}. Take $f$ as in (C4) and $g,y$ as in (C5), $B(x_1,x_2)=f(x_1)f(x_2)$,
and let $x\in\A^\kappa$ be the point with archimedean component $0$ and finite component $y$. Evaluation at $(x,x)$ and at $x$ are admissible
linear forms. By \cref{pf:sh-misc:main2} for $\omega^{(2)}$, the derivative at $0$ of $t\mapsto(\omega^{(2)}(\delta(t))(B\otimes(g\otimes g)))(x,x)$ is
\[
(8\pi i)^{-1}\sum_j\bigl(D^{(2)}_{(u_j|0)}D^{(2)}_{(v_j|0)}B+D^{(2)}_{(v_j|0)}D^{(2)}_{(u_j|0)}B\bigr)(0)+\sum_wc'_w\tr Z_w\cdot B(0),
\]
because $g(y)^2=1$. By (C3) and $f(0)=1$, $(D^{(2)}_{(p|0)}D^{(2)}_{(q|0)}B)(0)=(D_pD_qf)(0)$, and $B(0)=1$. On the other hand,
$B\otimes(g\otimes g)=(f\otimes g)\boxtimes(f\otimes g)$, so by the restriction property of the doubling data
$\omega^{(2)}(\delta(t))(B\otimes(g\otimes g))(x,x)=(\omega(\gamma(t))(f\otimes g))(x)\cdot(f\otimes g)(x)$, and $(f\otimes g)(x)=1$; by
\cref{pf:sh-misc:main2} for $\omega$, the derivative of this at $0$ is $(8\pi i)^{-1}\sum_j(D_{u_j}D_{v_j}f+D_{v_j}D_{u_j}f)(0)+\sum_wc_w\tr Z_w$.
Comparing the two values of the same derivative and using $\tr Z_w=\tr Y_w$ gives Step 1.

\emph{Step 2: $\sum_wc'_w\tr Y_w=\sum_w\frac{m_w}2\tr Y_w$.} Now let $\delta(t)=\mathrm{diag}(\gamma(t),\gamma(t))$ (C1), with
$\tau_w(\delta(t))=\exp(t\,\mathrm{diag}(Z_w,Z_w))$ and $\tr\mathrm{diag}(Z_w,Z_w)=2\tr Y_w$. Let $h\in\A_F$ be the archimedean adele $\frac12$. The
family consisting of $\bigl((hu_j|hu_j),(v_j|-v_j)\bigr)$ and $\bigl((hu_j|-hu_j),(v_j|v_j)\bigr)$ is a decomposition of $\mathrm{diag}(Z,Z)$: for
$x=(x_1|x_2)$, (C2) gives $\langle(v_j|-v_j),x\rangle^{(2)}=\langle v_j,x_1\rangle+\langle v_j,x_2\rangle$ and
$\langle(v_j|v_j),x\rangle^{(2)}=\langle v_j,x_1\rangle-\langle v_j,x_2\rangle$, so the sum of the corresponding terms is
$\sum_j(\langle v_j,x_1\rangle u_j\,|\,\langle v_j,x_2\rangle u_j)=(x_1{}^tZ\,|\,x_2{}^tZ)$.

Let $\ell,B,G$ be as in (C6); as $\ell$ is admissible, $\ell(B'\otimes G)=\Lambda(B')$ for a continuous linear form $\Lambda$ on archimedean
Schwartz functions, and $\Lambda(B)\ne0$. (a) \emph{$\Lambda(D^{(2)}_{(z|z)}B')=0$ for every archimedean $z$ and every $B'$:} the function
$t\mapsto\ell(\rho^{(2)}(t(z|z))(B'\otimes G))$ is constant by the invariance of $\ell$, and its derivative at $0$ is $\Lambda(D^{(2)}_{(z|z)}B')$.
(b) \emph{The Siegel rule.} The element $\delta(t)$ maps the diagonal ${1\choose1}$ to ${1\choose1}\gamma(t)$; the norm of $\det\gamma(t)$ is
$1$, because $\det\gamma(t)-1$ is archimedean and $|\tau_w(\det\gamma(t))|=|\exp(t\tr Y_w)|=1$ (\lref{PositivityTheory.RL.Up.norm_eq_one}),
and $\chi(\det\gamma(t))=\prod_w\exp(t\tr Y_w)^{m_w}$ by the infinity type of $\chi$. Hence the Siegel rule gives
\[
\ell\bigl(\omega^{(2)}(\delta(t))(B\otimes G)\bigr)=\exp\Bigl(t\sum_wm_w\tr Y_w\Bigr)\,\ell(B\otimes G),
\]
whose derivative at $0$ is $\bigl(\sum_wm_w\tr Y_w\bigr)\Lambda(B)$. (c) By \cref{pf:sh-misc:main2} for $\omega^{(2)}$ and the decomposition above,
the same derivative equals $\Lambda(\Theta B)+2\bigl(\sum_wc'_w\tr Y_w\bigr)\Lambda(B)$, where $8\pi i\,\Theta$ is the sum over $j$ of
\[
D^{(2)}_{a}D^{(2)}_{b}+D^{(2)}_{b}D^{(2)}_{a}+D^{(2)}_{a'}D^{(2)}_{b'}+D^{(2)}_{b'}D^{(2)}_{a'},
\]
with $a=(hu_j|hu_j)$, $b=(v_j|-v_j)$, $a'=(hu_j|-hu_j)$ and $b'=(v_j|v_j)$.
By (a), $\Lambda$ kills $D^{(2)}_aD^{(2)}_bB$ and $D^{(2)}_{b'}D^{(2)}_{a'}B$, as $a$ and $b'$ are diagonal. By the commutation relation (C7),
$\Lambda(D^{(2)}_bD^{(2)}_aB)=\Lambda(D^{(2)}_aD^{(2)}_bB)+2\pi i\langle b,a\rangle^{(2)}\Lambda(B)=2\pi i\langle b,a\rangle^{(2)}\Lambda(B)$ and
similarly $\Lambda(D^{(2)}_{a'}D^{(2)}_{b'}B)=\Lambda(D^{(2)}_{b'}D^{(2)}_{a'}B)+2\pi i\langle a',b'\rangle^{(2)}\Lambda(B)
=2\pi i\langle a',b'\rangle^{(2)}\Lambda(B)$. By (C2), $\langle b,a\rangle^{(2)}=2\langle v_j,hu_j\rangle$ and
$\langle a',b'\rangle^{(2)}=2\langle hu_j,v_j\rangle=-2\langle v_j,hu_j\rangle$, so the two terms cancel and $\Lambda(\Theta B)=0$. Comparing (b)
and (c) and dividing by $\Lambda(B)\ne0$ gives $2\sum_wc'_w\tr Y_w=\sum_wm_w\tr Y_w$.

\emph{Conclusion.} By Steps 1 and 2, \eqref{eq:pf:sh-misc:lieder} equals $\Phi\bigl(\sum_w(\shmiscR(Y_w)+\frac{m_w}2\tr Y_w)_wT\bigr)$, which is
the assertion of $\sfname{LieSecond}$.
\end{proof}

\begin{remark}\label{pf:sh-misc:lie2rem}
The proof of \cref{pf:sh-misc:main2} in the development combines $\sfname{LocalLift}$ with the fact that $\omega(g)$ and a local lift implementing
the same symplectic transformation in the Schr\"odinger model differ by a scalar factor, so that the derivative is $q(\rho(Z))f$ plus a
scalar $\zeta(Z)$; $\zeta$ is a real linear form on the Lie algebra of $\prod_w\Uni(K_w)$ vanishing on commutators
(\lref{PositivityTheory.ThKR.scalar_lie}), hence a combination $\sum_wc_w\tr Z_w$ of traces
(\lref{PositivityTheory.RL.Up.real_form_trace}), and $q(\rho(Z))$ is expressed through the operators $D_u$ by a decomposition of
$Z$ (\lref{PositivityTheory.ThKR.mid2}). These steps are not converted here.
\end{remark}

% ===========================================================================
\subsection{Howe's unramified generation}\label{sec:pf:sh-misc:howe}

This subsection proves the statement of the development which plays the role of Howe's unramified computation
\cite{Howe1979}: in an unramified local dual pair of unitary groups, in a lattice model of the Weil representation, the
vectors fixed by the integral points of one group are generated by the vacuum under the other group. It is used in
\cref{sec:pf:bundles,sec:pf:liftinj,sec:pf:open0,sec:pf:stageU} (through the separation of rows, \cref{prop:families:fs}).
The formal statement is axiomatic: the representation is any complex vector space with operators satisfying the Heisenberg
relations and the equivariance relations of a lattice model; no topology and no irreducibility are assumed. The proof is a
direct lattice-model argument: the space is graded by the classes of $M_{n'\times n}(L)$ modulo integral matrices, each class
contributing a line, and an induction on the index of the class reduces the statement to the existence, for every class
supporting a fixed vector, of an element of the span of the translates of the vacuum whose top part is concentrated on the
orbit of the class.


\paragraph{Setting.}
Let $L$ be a field of characteristic $0$, complete for a discrete valuation $v$ (written multiplicatively), with valuation ring
$\OO$ and maximal ideal $\mathfrak m$, and let $\varpi\in L$, an involutive ring automorphism $\sigma$ of $L$, $q\in\N$ and $\delta\in L$ be
given. We say that $(\varpi,\sigma,q)$ is \emph{unramified} if $\sigma^2=1$, $v\circ\sigma=v$, $\sigma\varpi=\varpi$, $\varpi$ is a uniformiser, $v(\sigma
x-x)=1$ for some $x\in\OO$, $q\ge2$ and $|\OO/\mathfrak m|=q^2$; then $L$ is an unramified quadratic extension of $F=L^\sigma$, whose
residue field has $q$ elements. For a matrix $x$ over $L$ put $x^*=\sigma(x)^t$. For $J\in\GL_m(L)$ let $\Uni(J)=\{g\in\GL_m(L):g^*Jg=J\}$.

Let $n,n'\ge0$, $J_S\in\GL_n(L)$, $J_T\in\GL_{n'}(L)$, an additive character $\psi\colon L\to\C^\times$ (no continuity is assumed), a complex
vector space $M$, a map $\rho\colon\shmiscX\to\End_\C(M)$ on $\shmiscX=M_{n'\times n}(L)$ and a homomorphism
$\omega\colon\Uni(J_S)\times\Uni(J_T)\to\End_\C(M)$. Put $\shmiscXO=M_{n'\times n}(\OO)$,
\[
H(x,y)=\delta\,\tr\bigl(J_S^{-1}y^*J_Tx\bigr),\qquad T(x,y)=H(x,y)+\sigma H(x,y)\qquad(x,y\in\shmiscX),
\]
$K_S=\Uni(J_S)\cap\GL_n(\OO)$ and $K_T=\Uni(J_T)\cap\GL_{n'}(\OO)$. A vector $s\in M$ is a \emph{vacuum vector} if $\rho(u)s=s$ for all
$u\in\shmiscXO$.

\begin{definition}[Local pairs]\label{pf:sh-misc:locpair}
$(\rho,\omega)$ is a \emph{local pair} for $(\varpi,\sigma,q,\delta,J_S,J_T,\psi)$ if:
(L1) $(\varpi,\sigma,q)$ is unramified and $q$ is odd;
(L2) $J_S\in\GL_n(\OO)$, $J_T\in\GL_{n'}(\OO)$, and $J_S^*=J_S$, $J_T^*=J_T$;
(L3) $v(\delta)=1$ and $\sigma\delta=-\delta$;
(L4) $\psi(a)=1$ for all $a\in\OO$ with $\sigma a=a$, and $\psi(\varpi^{-1}a)\ne1$ for some such $a$;
(L5) $\rho(0)=1$ and $\rho(x)\rho(y)=\psi\bigl(\tfrac12T(x,y)\bigr)\rho(x+y)$;
(L6) $M$ is spanned by the vectors $\rho(x)s$, $x\in\shmiscX$, $s$ a vacuum vector;
(L7) $\omega(g,h)\rho(x)=\rho(hxg^{-1})\omega(g,h)$ for $g\in\Uni(J_S)$, $h\in\Uni(J_T)$, $x\in\shmiscX$;
(L8) $\omega(g,h)s=s$ for $g\in K_S$, $h\in K_T$ and every vacuum vector $s$.
\formal{GlobalTheory.Hk3.LocPair}\formal{GlobalTheory.Hk3.LocVac}\formal{Formalism.UG.IsUnramInv}
\end{definition}

\begin{theorem}[Howe's unramified generation]\label{pf:sh-misc:howe}
Let $(\rho,\omega)$ be a local pair, and let $s\ne0$ be a vacuum vector such that the vectors $\rho(x)s$, $x\in\shmiscX$, span $M$. Then:
\begin{enumerate}[label=(\alph*)]
\item every $w\in M$ with $\omega(g,1)w=w$ for all $g\in K_S$ lies in the span of the vectors $\omega(1,h)s$, $h\in\Uni(J_T)$;
\item every $w\in M$ with $\omega(1,h)w=w$ for all $h\in K_T$ lies in the span of the vectors $\omega(g,1)s$, $g\in\Uni(J_S)$.
\end{enumerate}
\formallabel{GlobalTheory.Hk3.howeUnram_holds}\formallabel{GlobalTheory.Hk3.HoweDuality.fixed_mem_span}
\end{theorem}

\begin{remark}[Relation with the literature]\label{pf:sh-misc:howerem}
For the Weil representation of the unramified dual pair $(\Uni(J_S),\Uni(J_T))$, $p\ne2$, realised in a lattice model in which the
self-dual lattice is $\shmiscXO$ and $s$ is its vacuum (the characteristic function of the lattice), statement (b) is Howe's
theorem that the $K_T$-invariants are generated by the vacuum under $\Uni(J_S)$ \cite{Howe1979}; see
\cite[Ch.~5]{MoeglinVignerasWaldspurger1987} and \cite{Waldspurger1990} for proofs. The formal statement differs in that $M$ is an
arbitrary complex vector space with operators satisfying (L5)--(L8): there is no topology, $\psi$ is not assumed continuous, and
$M$ is not assumed irreducible (it is cyclic, generated by $s$ under the $\rho(x)$). It is proved in the development, as below, and
no published result is used. Statement (a) follows from (b) by exchanging the roles of the two groups.
\end{remark}

\paragraph{The lattice model.}
For a class $c\in\shmiscX/\shmiscXO$ let $V_c$ be the span of the vectors $\rho(x)s$, $x\in c$; for $x\in\shmiscX$ let
$\shmiscla(x)$ be the number of elements of $(x\OO^n+\OO^{n'})/\OO^{n'}$ (the \emph{index} of $x$). Throughout, $(\rho,\omega)$ is a local
pair and $s$ a vacuum vector.

\begin{lemma}[Lines of the lattice model]\label{pf:sh-misc:lattice}
\begin{enumerate}[label=(\alph*)]
\item For every vacuum vector $t$, $x\in\shmiscX$ and $u\in\shmiscXO$: $\rho(x+u)t=\psi(-\frac12T(x,u))\rho(x)t$ and
$\rho(u)\rho(x)t=\psi(T(u,x))\rho(x)t$. In particular $V_c=\C\rho(x)s$ for every $x\in c$.
\item $\rho(x)s\ne0$ for every $x$, if $s\ne0$.
\item The subspaces $V_c$, $c\in\shmiscX/\shmiscXO$, are independent (their sum is direct).
\item For $g\in K_S$, $h\in K_T$: $\omega(1,h)\rho(x)s=\rho(hx)s$ and $\omega(g,1)\rho(x)s=\rho(xg^{-1})s$.
\item If the $\rho(x)s$ span $M$, then every vacuum vector lies in $\C s$.
\end{enumerate}
\formallabel{GlobalTheory.Hk3.HoweDuality.rho_add_int}
\formallabel{GlobalTheory.Hk3.HoweDuality.rho_mem_of_mk_eq}
\formallabel{GlobalTheory.Hk3.locVac_span_classes_indep}
\formallabel{GlobalTheory.T.latticeModel_vacuumTranslate_classes_independent}
\formallabel{GlobalTheory.Hk3.HoweDuality.om_rho}
\formallabel{GlobalTheory.Hk3.locVac_mem_span_vac}
\formallabel{GlobalTheory.T.localLatticeVacuum_mem_span_fixed}
\end{lemma}

\begin{proof}
We first note $\sigma H(x,y)=-H(y,x)$: since $J_S,J_T$ are hermitian and $\sigma\delta=-\delta$, applying $\sigma$ to
$\delta\tr(J_S^{-1}y^*J_Tx)$ gives $-\delta\tr(x^*J_TyJ_S^{-1})=-H(y,x)$. Hence $T(x,y)=H(x,y)-H(y,x)=-T(y,x)$.

(a) By (L5), $\rho(x)t=\rho(x)\rho(u)t=\psi(\frac12T(x,u))\rho(x+u)t$, and $\psi(a)\psi(-a)=1$. Then
$\rho(u)\rho(x)t=\psi(\frac12T(u,x))\rho(u+x)t=\psi(\frac12T(u,x)-\frac12T(x,u))\rho(x)t=\psi(T(u,x))\rho(x)t$. The first formula shows
$\rho(y)s\in\C\rho(x)s$ whenever $y\equiv x\bmod\shmiscXO$.

(b) By (L5), $\rho(-x)\rho(x)s=\psi(\frac12T(-x,x))\rho(0)s$, a non-zero multiple of $s$.

(c) For $u\in\shmiscXO$ and $x\in c$ put $\chi_c(u)=\psi(T(u,x))$; by (a), $\rho(u)$ acts on $V_c$ by the scalar $\chi_c(u)$, which does not
depend on the choice of $x\in c$ (as $V_c$ is a line, or $0$). If $\chi_c=\chi_{c'}$ for $c=[x]$, $c'=[y]$, then $\psi(T(u,x-y))=1$ for all
$u\in\shmiscXO$, and the self-duality of $\shmiscXO$ for the pairing $\psi\circ T$ (a consequence of (L2)--(L4), as $L/F$ is unramified
and $\delta$ is a unit: \lref{GlobalTheory.Hk3.intMat_of_psi}) gives $x-y\in\shmiscXO$, i.e.\ $c=c'$. Thus the $V_c$ are eigenspaces of
the commuting family $(\rho(u))_{u\in\shmiscXO}$ for pairwise distinct systems of eigenvalues, and are therefore independent: a relation
$\sum_cw_c=0$ with $w_c\in V_c$ and minimal number of non-zero terms would, after applying $\rho(u)-\chi_{c_0}(u)$ with $\chi_{c_1}(u)\ne
\chi_{c_0}(u)$, give a shorter one.

(d) By (L7) and (L8), $\omega(1,h)\rho(x)s=\rho(hx)\omega(1,h)s=\rho(hx)s$, and similarly for $g$.

(e) Let $t$ be a vacuum vector and write $t=\sum_cw_c$ with $w_c\in V_c$ (possible by (c) and the spanning assumption). For
$u\in\shmiscXO$, $t=\rho(u)t=\sum_c\chi_c(u)w_c$, so $w_c=\chi_c(u)w_c$ by (c). If $w_c\ne0$, then $\chi_c=1=\chi_{[0]}$, so $c=[0]$ by the
argument in (c). Hence $t=w_{[0]}\in V_{[0]}=\C\rho(0)s=\C s$.
\end{proof}

\begin{lemma}[The index]\label{pf:sh-misc:latidx}
Let $(\varpi,\sigma,q)$ be unramified and $a,b\in\N$.
\begin{enumerate}[label=(\alph*)]
\item For $x\in M_{a\times b}(L)$ the index $\shmiscla(x)=|(x\OO^b+\OO^a)/\OO^a|$ is finite.
\item $\shmiscla(x+u)=\shmiscla(x)$ for $u\in M_{a\times b}(\OO)$, and $\shmiscla(hx)=\shmiscla(x)$ for $h\in\GL_a(\OO)$.
\item If every $u\in\OO^b$ with $xu\in\OO^a$ satisfies $yu\in\OO^a$, then $\shmiscla(y)\le\shmiscla(x)$.
\end{enumerate}
\formallabel{GlobalTheory.Hk3.HoweDuality.finite_range_latIdx}
\formallabel{GlobalTheory.Hk3.HoweDuality.latIdx_invariance}
\formallabel{GlobalTheory.T.matrixLatticeIndex_integral_translation_invariance}
\formallabel{GlobalTheory.Hk3.HoweDuality.latIdx_le_of_ker}
\formallabel{GlobalTheory.T.matrixLatticeIndex_le_of_integral_kernel}
\end{lemma}

\begin{proof}
(a) The entries of $x$ lie in $\varpi^{-k}\OO$ for some $k$, so $x\OO^b+\OO^a\subseteq\varpi^{-k}\OO^a$, and $\varpi^{-k}\OO^a/\OO^a$ is finite because
$\OO/\varpi\OO$ is finite (of order $q^2$). (b) $(x+u)\OO^b+\OO^a=x\OO^b+\OO^a$, and $h$ induces a bijection of $(x\OO^b+\OO^a)/\OO^a$ onto
$(hx\OO^b+\OO^a)/\OO^a$ since $h\OO^a=\OO^a$. (c) The map $u\mapsto xu\bmod\OO^a$ identifies $\OO^b/K_x$ with $(x\OO^b+\OO^a)/\OO^a$, where
$K_x=\{u\in\OO^b:xu\in\OO^a\}$; so $\shmiscla(x)=[\OO^b:K_x]$. The hypothesis is $K_x\subseteq K_y$, whence $[\OO^b:K_y]\le[\OO^b:K_x]$.
\end{proof}

\paragraph{Separation by height.}
The induction is carried out in the following abstract setting. Let a group $G$ act on a set $I$, let $M=\bigoplus_{i\in I}V_i$ be an
internal direct sum of subspaces of dimension at most $1$, and let $\pi\colon G\to\End_\C(M)$ be a homomorphism with $\pi(g)V_i\subseteq
V_{gi}$. For $m\in M$ let $m_i$ be its component in $V_i$; for $i\in I$ let $W_i=\sum_{g\in G}V_{gi}$ be the sum over the orbit of $i$, and for
$A\subseteq I$ let $W_A=\bigoplus_{j\in A}V_j$. Note that $(\pi(g)m)_{gk}=\pi(g)m_k$ for all $m\in M$, $g\in G$ and $k\in I$, because
$\pi(g)m_j\in V_{gj}$ and $gj\ne gk$ for $j\ne k$.

\begin{lemma}[Fixed vectors on one orbit]\label{pf:sh-misc:orbitline}
Let $i\in I$ and let $x,y\in W_i$ be $G$-fixed, with $y\ne0$. Then $x\in\C y$.
\formallabel{GlobalTheory.GradedFixed.mem_span_of_fixed}
\end{lemma}

\begin{proof}
As $y\ne0$ and $y\in W_i$, there is $k$ in the orbit of $i$ with $y_k\ne0$. Since $\dim V_k\le1$, $x_k=ay_k$ for some $a\in\C$. Then
$z=x-ay\in W_i$ is $G$-fixed and $z_k=0$. For $l=gk$ in the orbit, $z_l=(\pi(g)z)_{gk}=\pi(g)z_k=0$; for $l$ outside the orbit, $z_l=0$ as
$z\in W_i$. Hence $z=0$, i.e.\ $x=ay$.
\end{proof}

\begin{lemma}[Orbit-height separation]\label{pf:sh-misc:orbitheight}
Let $\mu\colon I\to\N$ be $G$-invariant and $S\subseteq M$ a subspace of $G$-fixed vectors. For $i\in I$ let
$W_{<i}=\sum_{\mu(j)<\mu(i)}V_j$. Assume: for every $i$ such that $W_i$ contains a non-zero $G$-fixed vector there are $v\in S$ and $u\in
W_{<i}$ with $v-u\in W_i$ and $v-u\ne0$. Then every $G$-fixed vector of $M$ lies in $S$.
\formallabel{GlobalTheory.fixed_mem_of_orbit_height}
\formallabel{GlobalTheory.T.invariant_mem_of_orbit_height_separation}
\end{lemma}

\begin{proof}
\emph{If $m$ is $G$-fixed and $A\subseteq I$ is $G$-stable, the projection $m_A=\sum_{j\in A}m_j$ of $m$ to $W_A$ is $G$-fixed}: $m=m_A+m_{I\smallsetminus A}$,
both summands are mapped into themselves by $\pi(g)$ (as $A$ and $I\smallsetminus A$ are $G$-stable), and comparing components of
$\pi(g)m=m$ gives $\pi(g)m_A=m_A$.

We show by strong induction on $\mu(i)$: \emph{every $G$-fixed $x\in W_i$ lies in $S$.} Let $x\ne0$ be $G$-fixed in $W_i$ and take $v,u$ as
in the hypothesis. The set $\{j:\mu(j)<\mu(i)\}$ is $G$-stable and disjoint from the orbit of $i$; since $v=(v-u)+u$ is $G$-fixed with
$v-u\in W_i$ and $u$ supported outside the orbit, the first step shows that $v-u$ and $u$ are $G$-fixed. Now $u$ is a finite sum of its
projections to the orbit sums $W_j$ with $\mu(j)<\mu(i)$, each $G$-fixed by the first step, hence in $S$ by induction; so $u\in S$ and
$v-u\in S$. By \cref{pf:sh-misc:orbitline}, $x$ is a multiple of the non-zero fixed vector $v-u$, so $x\in S$. Finally, an arbitrary $G$-fixed
$w\in M$ is the finite sum of its projections to the orbit sums meeting its support, each of which is $G$-fixed by the first step and lies in
$S$ by what was just shown.
\end{proof}

\paragraph{The induction step.}
We use the following results of the formal development, which are part of this proof but are not converted here.
For $x,y\in\shmiscX$ write $\mathrm{Gram}(x)=xJ_S^{-1}x^*\in M_{n'}(L)$ and $h_T(a,b)=a^*J_Tb$ for $a,b\in L^{n'}$; say that
$\mathrm{HomA}(x,y)$ holds if every $u\in\OO^n$ with $xu\in\OO^{n'}$ satisfies $yu\in\OO^{n'}$, and $yJ_S^{-1}x^*J_Tw\in\OO^{n'}$ for every
$w\in\OO^{n'}$; and that $\mathrm{GramEq}(x,y)$ holds if $c\bigl(h_T(ym,ym')-h_T(xm,xm')\bigr)\in\OO$ for all $m,m'\in\OO^n$ and $c\in\OO$ with
$x(cm),x(cm')\in\OO^{n'}$. A matrix $a\in\GL_n(L)$ is \emph{self-dual} if $a^*J_Sa$ and $a^{-1}J_S^{-1}(a^{-1})^*$ are integral.

\begin{enumerate}[label=(D\arabic*)]
\item\label{pf:sh-misc:D-cayley} (\lref{GlobalTheory.Hk3.CayleyPhase.cayley_phase},
\lref{GlobalTheory.T.cayleyTransform_trivial_phase_skewHermitian}) Let $x\in\shmiscX$, and assume $\psi(\frac12T(x,hx-x))=1$ for every $h\in
K_T$ with $hx-x\in\shmiscXO$. Then $\psi\bigl(\frac12(\delta\tr(C\,\mathrm{Gram}(x))+\sigma(\delta\tr(C\,\mathrm{Gram}(x))))\bigr)=1$ for every
$C\in M_{n'}(L)$ with $C^*=-C$, entries in $\mathfrak m$ and $Cx\in\shmiscXO$.
\item\label{pf:sh-misc:D-phases} (\lref{GlobalTheory.Hk3.HoweDuality.exists_int_gram_of_phases},
\lref{GlobalTheory.T.exists_integral_gram_of_phase_conditions}) If $x$ satisfies the conclusion of \ref{pf:sh-misc:D-cayley}, there is
$l\in\shmiscXO$ such that $\mathrm{Gram}(x-l)$ is integral.
\item\label{pf:sh-misc:D-sum} (\lref{GlobalTheory.Hk3.HoweG.gramClass_sum_mem_span_selfDual_vac},
\lref{GlobalTheory.T.gram_class_translate_sum_mem_span_selfDual_lattice_vacua}) If $\mathrm{Gram}(x)$ is integral, there are a finite set
$C\ni x$ of elements of $\shmiscX$, pairwise incongruent modulo $\shmiscXO$, and non-zero $c(y)\in\C$ ($y\in C$), such that $\mathrm{HomA}(x,y)$ and
$\mathrm{GramEq}(x,y)$ hold for $y\in C$ and $\sum_{y\in C}c(y)\rho(y)s$ lies in the span of the vectors $w$ for which there is a self-dual
$a$ with $\rho(ua^{-1})w=w$ for all $u\in\shmiscXO$.
\item\label{pf:sh-misc:D-witt} (\lref{GlobalTheory.Hk3.HoweG.exists_intUnitary_of_gramClass},
\lref{GlobalTheory.T.exists_intUnitary_of_gramClass_alt}) If $\mathrm{Gram}(x)$ is integral, $\mathrm{HomA}(x,y)$, $\mathrm{GramEq}(x,y)$
and $\shmiscla(y)=\shmiscla(x)$, then $y-hx\in\shmiscXO$ for some $h\in K_T$.
\end{enumerate}

\begin{lemma}[Self-dual matrices]\label{pf:sh-misc:selfdual}
Assume \textup{(L1)} and \textup{(L2)}. Every self-dual $a\in\GL_n(L)$ is of the form $a=gk$ with $g\in\Uni(J_S)$ and $k\in\GL_n(\OO)$.
\formallabel{GlobalTheory.Hk3.HoweG.selfDual_eq_unitary_mul_int}
\formallabel{GlobalTheory.T.selfDual_lattice_eq_unitary_mul_integralGL}
\end{lemma}

\begin{proof}
We use that every hermitian $J\in\GL_n(\OO)$ (i.e.\ $J^*=J$) satisfies $P^*JP=J_0$ for some $P\in\GL_n(\OO)$, where $J_0$ is a fixed
antidiagonal matrix depending only on $n$ (\lref{GlobalTheory.NormalHyp.frame}; this is the uniqueness of unimodular hermitian lattices
over the unramified extension $L/F$ with $q$ odd, compare \cite{Jacobowitz1962}). The matrix $G=a^*J_Sa$ is hermitian, since $J_S$ is,
and $G$ and $G^{-1}=a^{-1}J_S^{-1}(a^{-1})^*$ are integral because $a$ is self-dual; so $G\in\GL_n(\OO)$. Choose $P,Q\in\GL_n(\OO)$ with
$P^*GP=J_0=Q^*J_SQ$, and put $g=aPQ^{-1}$ and $k=QP^{-1}\in\GL_n(\OO)$. Then $a=gk$ and
$g^*J_Sg=(Q^{-1})^*\,P^*a^*J_SaP\,Q^{-1}=(Q^{-1})^*J_0Q^{-1}=J_S$, i.e.\ $g\in\Uni(J_S)$.
\end{proof}

\begin{lemma}[Classes supporting fixed vectors]\label{pf:sh-misc:classfixed}
Let $x\in\shmiscX$ and assume that the span of the $\rho(hx)s$, $h\in K_T$, contains a non-zero vector fixed by all $\omega(1,h)$,
$h\in K_T$. Then $\rho(x)s\ne0$, and $\rho(hx)s=\rho(x)s$ for every $h\in K_T$ with $hx-x\in\shmiscXO$. Consequently
$\psi(\frac12T(x,hx-x))=1$ for every such $h$.
\formallabel{GlobalTheory.Hk3.class_fixed_phase}
\formallabel{GlobalTheory.T.heisenbergTranslate_fixed_integralUnitary}
\end{lemma}

\begin{proof}
Let $w\ne0$ be such a fixed vector; it lies in $W=\sum_{h\in K_T}V_{[hx]}$. By \cref{pf:sh-misc:lattice}(d), $\omega(1,h)$ maps $V_c$ into
$V_{hc}$ (where $h[x]=[hx]$) and is invertible. If the component $w_{[x]}$ of $w$ in $V_{[x]}$ vanished, then so would all components
$w_{[hx]}=(\omega(1,h)w)_{[hx]}=\omega(1,h)w_{[x]}$ (as in the setting of \cref{pf:sh-misc:orbitline}), and $w=0$. Hence $w_{[x]}\ne0$, so $\rho(x)s\ne0$ and
$w_{[x]}=a\rho(x)s$ with $a\ne0$. If $h\in K_T$ and $hx-x\in\shmiscXO$, then $h$ fixes the class $[x]$, so comparing the $[x]$-components of
$\omega(1,h)w=w$ gives $a\rho(hx)s=\omega(1,h)w_{[x]}=w_{[x]}=a\rho(x)s$. For the last assertion, (L5) and the vacuum property give
$\rho(x)s=\rho(x)\rho(hx-x)s=\psi(\frac12T(x,hx-x))\rho(hx)s=\psi(\frac12T(x,hx-x))\rho(x)s$, and $\rho(x)s\ne0$.
\end{proof}

\begin{lemma}[Integral Gram matrices]\label{pf:sh-misc:intgram}
Let $s$ be a vacuum vector and $x\in\shmiscX$ such that the span of the $\rho(hx)s$, $h\in K_T$, contains a non-zero vector fixed by all
$\omega(1,h)$, $h\in K_T$. Then there is $l\in\shmiscXO$ such that $\mathrm{Gram}(x-l)$ is integral.
\formallabel{GlobalTheory.Hk3.HoweDuality.exists_int_gram_of_fixed}
\formallabel{GlobalTheory.T.exists_int_gram_of_fixed_alt}
\end{lemma}

\begin{proof}
By \cref{pf:sh-misc:classfixed}, $\psi(\frac12T(x,hx-x))=1$ for all $h\in K_T$ with $hx-x\in\shmiscXO$. By \ref{pf:sh-misc:D-cayley}, $x$
satisfies the phase conditions for the skew-hermitian matrices $C$, and \ref{pf:sh-misc:D-phases} gives $l$.
\end{proof}

\begin{proposition}[Translates of the vacuum with top part on an orbit]\label{pf:sh-misc:gramorbit}
Let $s\ne0$ be a vacuum vector such that the $\rho(x)s$ span $M$, and $S$ the span of the $\omega(g,1)s$, $g\in\Uni(J_S)$. Let $x\in\shmiscX$ with
$\mathrm{Gram}(x)$ integral. Then there are $v\in S$ and $u$ in the span of the $\rho(y)s$ with $\shmiscla(y)<\shmiscla(x)$ such that $v-u$ is a
non-zero element of the span of the $\rho(hx)s$, $h\in K_T$.
\formallabel{GlobalTheory.Hk3.HoweG.gramClass_orbit_mod_lower}
\formallabel{GlobalTheory.T.exists_conjugateModel_extremal_dual_parameter}
\end{proposition}

\begin{proof}
Take $C$ and $c$ as in \ref{pf:sh-misc:D-sum} and put $v=\sum_{y\in C}c(y)\rho(y)s$.

\emph{$v\in S$.} Let $w$ and a self-dual $a$ be as in \ref{pf:sh-misc:D-sum}, and write $a=gk$ by \cref{pf:sh-misc:selfdual}. For $u\in\shmiscXO$ we have
$uk\in\shmiscXO$ and $(uk)a^{-1}=ug^{-1}$, so $\rho(ug^{-1})w=w$. By (L7),
$\rho(u)\omega(g^{-1},1)w=\omega(g^{-1},1)\rho(ug^{-1})w=\omega(g^{-1},1)w$, so $\omega(g^{-1},1)w$ is a vacuum vector, hence a multiple of $s$
(\cref{pf:sh-misc:lattice}(e)), and $w\in\C\omega(g,1)s\subseteq S$. Thus $v\in S$.

\emph{The lower part.} By $\mathrm{HomA}(x,y)$ and \cref{pf:sh-misc:latidx}(c), $\shmiscla(y)\le\shmiscla(x)$ for all $y\in C$. Let $u$ be the sum of
the terms $c(y)\rho(y)s$ with $\shmiscla(y)\ne\shmiscla(x)$, i.e.\ $\shmiscla(y)<\shmiscla(x)$. Then $v-u=\sum_{y\in P}c(y)\rho(y)s$ with
$P=\{y\in C:\shmiscla(y)=\shmiscla(x)\}\ni x$.

\emph{$v-u$ lies on the orbit.} For $y\in P$, \ref{pf:sh-misc:D-witt} gives $h\in K_T$ with $y-hx\in\shmiscXO$, so $\rho(y)s\in\C\rho(hx)s$ by
\cref{pf:sh-misc:lattice}(a).

\emph{$v-u\ne0$.} The elements of $P\smallsetminus\{x\}$ are not congruent to $x$ modulo $\shmiscXO$, so $\sum_{y\in P\smallsetminus\{x\}}c(y)\rho(y)s$ lies
in $\sum_{c'\ne[x]}V_{c'}$. If $v-u=0$, the independence of the $V_c$ (\cref{pf:sh-misc:lattice}(c)) would give $c(x)\rho(x)s=0$, which is
impossible because $c(x)\ne0$ and $\rho(x)s\ne0$ (\cref{pf:sh-misc:lattice}(b)).
\end{proof}

\begin{proposition}[The induction step]\label{pf:sh-misc:howestep}
Let $s\ne0$ be a vacuum vector such that the $\rho(x)s$ span $M$, and $S$ the span of the $\omega(g,1)s$, $g\in\Uni(J_S)$. Let $x\in\shmiscX$
such that the span of the $\rho(hx)s$, $h\in K_T$, contains a non-zero vector fixed by $\omega(1,K_T)$. Then there are $v\in S$ and $u$ in the
span of the $\rho(y)s$ with $\shmiscla(y)<\shmiscla(x)$ such that $v-u$ is a non-zero element of the span of the $\rho(hx)s$, $h\in K_T$.
\formallabel{GlobalTheory.Hk3.HoweDuality.step}
\end{proposition}

\begin{proof}
By \cref{pf:sh-misc:intgram} there is $l\in\shmiscXO$ with $\mathrm{Gram}(x')$ integral, $x'=x-l$. By \cref{pf:sh-misc:latidx}(b),
$\shmiscla(x')=\shmiscla(x)$. \Cref{pf:sh-misc:gramorbit} applied to $x'$ gives $v\in S$ and $u$ with $v-u\ne0$ in the span of the
$\rho(hx')s$; since $hx'-hx=-hl\in\shmiscXO$, $\rho(hx')s\in\C\rho(hx)s$ (\cref{pf:sh-misc:lattice}(a)), so $v-u$ lies in the span of the
$\rho(hx)s$.
\end{proof}

\begin{proof}[Proof of \cref{pf:sh-misc:howe}]
(b) Let $G=K_T$ act on $I=\shmiscX/\shmiscXO$ by $h[x]=[hx]$ (well defined since $h\shmiscXO=\shmiscXO$), let $V_c$ be as above,
$\pi(h)=\omega(1,h)$, $\mu([x])=\shmiscla(x)$ (well defined and $K_T$-invariant by \cref{pf:sh-misc:latidx}(b)) and $S$ the span of the
$\omega(g,1)s$. The vectors of $S$ are $K_T$-fixed: $\omega(1,h)\omega(g,1)s=\omega(g,1)\omega(1,h)s=\omega(g,1)s$ by (L8). By
\cref{pf:sh-misc:lattice}, $M=\bigoplus_cV_c$ with $\dim V_c\le1$ and $\pi(h)V_c\subseteq V_{hc}$. For $c=[x]$, $W_c$ is the span of the
$\rho(hx)s$, $h\in K_T$, and $W_{<c}$ is the span of the $\rho(y)s$ with $\shmiscla(y)<\shmiscla(x)$. So the hypothesis of
\cref{pf:sh-misc:orbitheight} is \cref{pf:sh-misc:howestep}, and \cref{pf:sh-misc:orbitheight} gives (b).

(a) Let $\rho'(X)=\rho(J_T^{-1}X^*J_S)$ for $X\in M_{n\times n'}(L)$ and $\omega'(h,g)=\omega(g,h)$. Then $(\rho',\omega')$ is a local pair for
$(\varpi,\sigma,q,-\delta,J_T,J_S,\psi)$, the vacuum vectors of $\rho$ are vacuum vectors of $\rho'$, and the $\rho'(X)s$ span $M$ if the $\rho(x)s$
do (\lref{GlobalTheory.Hk3.HoweSwap.locPair_swap}: the map
$X\mapsto J_T^{-1}X^*J_S$ is a bijection $M_{n\times n'}(L)\to\shmiscX$ preserving integrality, and it transforms the cocycle of (L5) for
$(-\delta,J_T,J_S)$ into that for $(\delta,J_S,J_T)$). Statement (b) for $(\rho',\omega')$ is statement (a) for $(\rho,\omega)$.
\end{proof}

% ===========================================================================
\subsection{Primes represented by norm forms in short intervals}\label{sec:pf:sh-misc:shortnorm}

This subsection proves the lower bound for primes of the form $\pi\bar\pi$, $\pi\equiv1\bmod M$, in an imaginary quadratic field,
in short intervals, uniformly in the discriminant and the modulus. It is used in \cref{sec:pf:stageU} and, through
\lref{UpperTheory.T3.normPrimesN}, in \cref{sec:pf:top-supply}. The formal proof is axiomatic: it introduces
\emph{$L$-families} (a finite abelian group with arithmetic functions satisfying a lattice-point count, a theta bound, a
convolution identity and a positivity property), proves a short-interval lower bound for every $L$-family by the classical
analytic method (Mellin continuation, zero-free region, at most one exceptional zero, an effective bound for it, explicit formula),
and constructs an $L$-family for the ray classes of $\Q(\sqrt{-d})$ modulo $M$. We give the definitions, the deduction of the
theorem, and the final step of the analytic argument; the analytic inputs and the construction are stated precisely and cited.

\paragraph{The statement.}
For a prime $d\equiv3\bmod4$ and integers $M\ge1$ and $p$, we say that $p$ is a \emph{norm-one prime for $(d,M)$} if $p$ is prime and there are
$A,B\in\Z$ with $A^2+dB^2=4p$, $M\mid B$ and $2M\mid A-2-B$. For reals $X\le Y$ let $N_{d,M}(X,Y)$ be the number of norm-one primes $p$ for
$(d,M)$ with $\lfloor X\rfloor<p\le\lfloor Y\rfloor$.

\begin{theorem}[Norm-one primes in short intervals]\label{pf:sh-misc:shortnorm}
There are $A,c,c'>0$ such that for every prime $d\equiv3\bmod4$ with $d\ge100$, every integer $M\ge3$ and all reals $X\ge1$ and $h$ with
\[
\log(4dM^2)\le c\sqrt{\log X}\qquad\text{and}\qquad e^{-c\sqrt{\log X}}\le h\le1,
\]
one has
\[
N_{d,M}\bigl(X,X(1+h)\bigr)\ \ge\ \frac{c'\,hX}{(4dM^2)^A\log X}.
\]
\formallabel{UpperTheory.T3.shortNorm}
\formallabel{GlobalTheory.T.normPrimesImagQuad_short_interval_lower_bound}
\formallabel{GlobalTheory.Hk.normPrimes_short_interval_lower}
\end{theorem}

\begin{remark}\label{pf:sh-misc:shortnormrem}
Let $K=\Q(\sqrt{-d})$, with $\OO_K=\Z[\omega]$, $\omega=\frac{1+\sqrt{-d}}2$. If $A^2+dB^2=4p$, then $A\equiv B\bmod2$ and $\pi=\frac{A+B\sqrt{-d}}2=
\frac{A-B}2+B\omega\in\OO_K$ has norm $p$; and $\pi-1=\frac{A-B-2}2+B\omega\in M\OO_K$ if and only if $M\mid B$ and $2M\mid A-B-2$. So $p$ is a
norm-one prime for $(d,M)$ if and only if $p=\pi\bar\pi$ for some $\pi\in\OO_K$ with $\pi\equiv1\bmod M\OO_K$, which is how the statement is
used in \cref{sec:pf:stageU}. The statement is of the type of the prime number theorem for ray classes in short intervals with the
classical (de la Vall\'ee Poussin) error term, with a polynomial loss in the conductor which accounts for a possible exceptional zero
(compare \cite[Ch.~5]{IwaniecKowalski2004} and \cite{Hecke1920}); the formal development proves it as follows, without citing
the literature.
\end{remark}

\paragraph{$L$-families.}
\begin{definition}[$L$-families]\label{pf:sh-misc:lfam}
An \emph{$L$-family} $\mathcal F$ consists of a finite abelian group $G$, reals $B$ and $\kappa$, and functions $a,\lambda\colon G\times\N\to\R$
such that:
\begin{enumerate}[label=(F\arabic*)]
\item $B\ge4$, $|G|\le B$ and $1/B\le\kappa\le4$;
\item $a(g,n)\ge0$, $a(g,0)=0$, $a(0,1)=1$ and $a(g,1)=0$ for $g\ne0$;
\item (count) $\bigl|\sum_{1\le n\le x}a(g,n)-\kappa x\bigr|\le10\sqrt x$ for all $g$ and $x\ge1$;
\item (theta) $\bigl|\sum_{n\ge0}a(g,n)e^{-tn}-\kappa/t\bigr|\le B^{10}t$ for all $g$ and $0<t\le1$;
\item $\lambda(g,n)\ge0$ and $\sum_g\lambda(g,n)\le2\Lambda(n)$, where $\Lambda$ is the von Mangoldt function;
\item (convolution) $a(g,n)\log n=\sum_{h\in G}\sum_{n_1n_2=n}\lambda(h,n_1)\,a(g-h,n_2)$ for all $g$ and $n$;
\item (positivity) for every character $\chi$ of $G$ with $\chi^2=1$ and every $n$, the coefficient
$c_\chi(n)=\sum_{n_1n_2=n}\bigl(\sum_ga(g,n_1)\bigr)\bigl(\sum_g\Re\chi(g)\,a(g,n_2)\bigr)$ is $\ge0$, and $\sum_ga(g,n)\le c_\chi(n^2)$.
\end{enumerate}
\formal{GlobalTheory.Hk.LFam}
\end{definition}

For a character $\chi$ of $G$ put $a_\chi(n)=\sum_g\chi(g)a(g,n)$, $\lambda_\chi(n)=\sum_g\chi(g)\lambda(g,n)$, $\psi_\chi(u)=\sum_{1\le n\le
u}\lambda_\chi(n)$, $r_\chi=\kappa|G|$ if $\chi=1$ and $r_\chi=0$ otherwise,
\[
m_\chi(s)=\int_1^\infty\Bigl(\sum_{1\le n\le x}a_\chi(n)-r_\chi x\Bigr)x^{-s-1}\,dx,\qquad
F_\chi(s)=\begin{cases}r_\chi s+(s-1)s\,m_\chi(s)&\chi=1,\\ s\,m_\chi(s)&\chi\ne1,\end{cases}
\]
and $\ell(t)=\log\bigl(B(|t|+2)\bigr)$. By (F3), $m_\chi$ is holomorphic for $\Re s>\frac12$; for $\Re s>1$, $F_\chi(s)$ equals
$\sum_na_\chi(n)n^{-s}$, multiplied by $s-1$ if $\chi=1$. So $F_\chi$ is the continuation to $\Re s>\frac12$ of the $L$-series of $\chi$ (with the
pole removed). For $c>0$, a finite set $Z\subset\R$ is an \emph{exceptional set} for $(\chi,c)$ if $|Z|\le1$; every $\beta\in Z$ satisfies
$1-c/\ell(0)\le\beta<1$, $F_\chi(\beta)=0$ and $F'_\chi(\beta)\ne0$; and every zero $s$ of $F_\chi$ with $\Re s\ge1-c/\ell(\Im s)$ is real and lies in
$Z$.

We use the following analytic results of the formal development (they are part of this proof and are not converted here).
\begin{enumerate}[label=(Z\arabic*)]
\item\label{pf:sh-misc:Z-zf} (Zero-free region; \lref{GlobalTheory.T.rayClassLFunction_zero_free_region}) There is $c_4>0$ such that
for every $L$-family, every $\chi$ and every zero $s$ of $F_\chi$ with $\Re s\ge\frac34$ and ($\chi^2\ne1$ or $\Im s\ne0$),
$\Re s<1-c_4/\ell(\Im s)$.
\item\label{pf:sh-misc:Z-exc} (At most one exceptional zero;
\lref{GlobalTheory.T.rayClassLFunction_unique_simple_exceptional_zero}) There is $c_5>0$ such that for every $L$-family, all characters
$\chi,\chi'$ with $\chi^2=\chi'^2=1$ and all real zeros $\beta$ of $F_\chi$ and $\beta'$ of $F_{\chi'}$ with $\beta,\beta'\ge\max(\frac34,1-c_5/\ell(0))$:
$\chi=\chi'$, $\beta=\beta'$, and $F'_\chi(\beta)\ne0$.
\item\label{pf:sh-misc:Z-far} (Effective bound for the exceptional zero;
\lref{GlobalTheory.T.exists_rayClassLFunction_real_zero_far_from_one}) There are $A_0,c_9>0$ such that for every $L$-family, every $\chi$
with $\chi^2=1$ and every real zero $\beta\ge\frac34$ of $F_\chi$: $1-\beta\ge c_9/B^{A_0}$.
\item\label{pf:sh-misc:Z-psi} (Explicit formula; \lref{GlobalTheory.T.rayClassFamily_psi_with_exceptional_zero}) For every
$0<c_0\le\frac1{20}$ there are $c_7,C>0$ such that for every $L$-family, every $\chi$, every exceptional set $Z$ for $(\chi,c_0)$ and every $u\ge2$
with $\log B\le\sqrt{\log u}$,
\[
\Bigl|\psi_\chi(u)-[\chi=1]\,u+\sum_{\beta\in Z}\frac{u^\beta}\beta\Bigr|\ \le\ C\,u\,e^{-c_7\sqrt{\log u}}.
\]
\end{enumerate}

\begin{lemma}[Exceptional zeros]\label{pf:sh-misc:zeros}
There are $0<c_0\le\frac1{20}$, $A_0>0$ and $0<c_9'\le1$ such that for every $L$-family there are exceptional sets $Z(\chi)$ for $(\chi,c_0)$,
one for each character $\chi$, with: if $\beta\in Z(\chi)$ and $\beta'\in Z(\chi')$, then $\chi=\chi'$ and $\beta=\beta'$; and every $\beta\in Z(\chi)$
satisfies $\beta\ge\frac34$ and $1-\beta\ge c_9'/B^{A_0}$.
\formallabel{GlobalTheory.Hk.ShortInt.zeros}
\end{lemma}

\begin{proof}
Let $c_4,c_5,(A_0,c_9)$ be as in \ref{pf:sh-misc:Z-zf}--\ref{pf:sh-misc:Z-far}, $c_0=\min(\frac1{20},c_4,c_5)$ and $c_9'=\min(c_9,1)$. Note
$\ell(t)\ge\log8>1$ as $B\ge4$. \emph{(i) If $F_\chi(s)=0$ and $\Re s\ge1-c_0/\ell(\Im s)$, then $\Re s\ge\frac34$, $\chi^2=1$ and $\Im s=0$}: since
$c_0/\ell\le\frac1{20}$, $\Re s\ge\frac{19}{20}$; if $\chi^2\ne1$ or $\Im s\ne0$, \ref{pf:sh-misc:Z-zf} gives $\Re s<1-c_4/\ell(\Im s)\le1-c_0/\ell(\Im s)$, a
contradiction. \emph{(ii)} By (i), \ref{pf:sh-misc:Z-exc} applies to any two real zeros $\beta,\beta'\ge1-c_0/\ell(0)$ (of $F_\chi$, $F_{\chi'}$),
since $c_0\le c_5$: $\chi=\chi'$, $\beta=\beta'$ and $F'_\chi(\beta)\ne0$; and by \ref{pf:sh-misc:Z-far}, $1-\beta\ge c_9/B^{A_0}\ge c'_9/B^{A_0}>0$.
Let $Z(\chi)=\{\beta\}$ if $F_\chi$ has a real zero $\beta\ge1-c_0/\ell(0)$ (it is unique by (ii)) and $Z(\chi)=\varnothing$ otherwise. Then $|Z(\chi)|\le1$;
an element $\beta$ satisfies $\beta<1$, $F_\chi(\beta)=0$ and $F'_\chi(\beta)\ne0$ by (ii); and a zero $s$ of $F_\chi$ with $\Re s\ge1-c_0/\ell(\Im s)$ is real by
(i), hence equal to $\beta$. The remaining assertions are (i) and (ii).
\end{proof}

\begin{theorem}[Short intervals in $L$-families]\label{pf:sh-misc:thmA}
There are $A,c,c'>0$ such that for every $L$-family and all reals $X\ge1$ and $h$ with $\log B\le c\sqrt{\log X}$ and
$e^{-c\sqrt{\log X}}\le h\le1$,
\[
\sum_{\lfloor X\rfloor<n\le\lfloor X(1+h)\rfloor}\lambda(0,n)\ \ge\ \frac{c'\,hX}{B^A}.
\]
\formallabel{GlobalTheory.Hk.ShortInt.trivialClass_shortInterval}
\formallabel{GlobalTheory.T.exists_lam_zero_sum_shortInterval_ge}
\end{theorem}

\begin{proof}
Let $c_0,A_0,c_9'$ be as in \cref{pf:sh-misc:zeros} and $c_7,C$ as in \ref{pf:sh-misc:Z-psi} for this $c_0$. Put $A=A_0+1$,
$c=\min\bigl(1,\frac{c_7}{2(A_0+2)},\frac{c_7c'_9}{24C}\bigr)$ and $c'=c'_9/4$. Let an $L$-family and $X,h$ be given, and let $Z(\chi)$ be as in
\cref{pf:sh-misc:zeros}.

\emph{Ranges.} As $B\ge4$, $\log B\ge1$; so $\sqrt{\log X}\ge(\log B)/c\ge\log B\ge1$, whence $\log X\ge1$, $X\ge e$, and
$X\le X(1+h)\le2X$. Put $L=\sqrt{\log X}$ and $E=e^{-c_7L}$; then $\log B\le\sqrt{\log u}$ for $u\in\{X,X(1+h)\}$, and
$e^{-c_7\sqrt{\log(X(1+h))}}\le E$.

\emph{Orthogonality.} For $u\ge0$, $\sum_\chi\psi_\chi(u)=|G|\sum_{1\le n\le u}\lambda(0,n)$, by the orthogonality of the characters of $G$. Summing
\ref{pf:sh-misc:Z-psi} over the $|G|\le B$ characters gives, for $u=X$ and $u=X(1+h)$,
\[
\Bigl||G|\sum_{1\le n\le u}\lambda(0,n)-u+\sum_\chi\sum_{\beta\in Z(\chi)}\frac{u^\beta}\beta\Bigr|\le B\cdot Cu\,E\le2BCXE .
\]
Let $S=\sum_{\lfloor X\rfloor<n\le\lfloor X(1+h)\rfloor}\lambda(0,n)\ge0$. Subtracting the two estimates,
\[
|G|\,S\ \ge\ hX-\sum_\chi\sum_{\beta\in Z(\chi)}\Bigl(\frac{(X(1+h))^\beta}\beta-\frac{X^\beta}\beta\Bigr)-3BCXE .
\]

\emph{The exceptional term.} By \cref{pf:sh-misc:zeros} the double sum has at most one term. Let $\beta$ be its exceptional zero, so $\beta\ge\frac34$
and $1-\beta\ge\delta:=c_9'/B^{A_0}$, with $0<\delta\le1$. By Bernoulli's inequality $(1+h)^\beta\le1+\beta h$, and
$X^\beta=Xe^{-(1-\beta)\log X}\le Xe^{-\delta}\le X(1-\delta/2)$ (as $\log X\ge1$ and $e^{-\delta}\le1-\delta/2$ for $0<\delta\le1$). Hence
$\frac{(X(1+h))^\beta-X^\beta}\beta=\frac{X^\beta((1+h)^\beta-1)}\beta\le X^\beta h\le hX(1-\delta/2)$. The same bound holds trivially if there is no
exceptional zero.

\emph{The error term.} We claim $3BCE\le h\delta/4$. Indeed $B\le e^{cL}$, $B^{A_0}\le e^{cLA_0}$, and, since $cL\ge1$ and $c\le c_7c'_9/(24C)$,
$12C\le12C\,cL\le\frac12c_7c'_9L\le c'_9e^{c_7L/2}$. As
$(A_0+2)c\le c_7/2$,
\[
3BCE\cdot4B^{A_0}=B\,B^{A_0}(12C)E\le e^{cL+cLA_0+c_7L/2-c_7L}\,c'_9\le e^{-cL}c'_9\le h\,c'_9 .
\]

\emph{Conclusion.} Combining, $|G|S\ge hX-hX(1-\delta/2)-hX\delta/4=hX\delta/4$, and $|G|\le B$ gives
$S\ge\frac{c'_9}4\,hX/B^{A_0+1}=c'hX/B^A$.
\end{proof}

\begin{theorem}[An $L$-family for ray classes]\label{pf:sh-misc:lfamray}
For every prime $d\equiv3\bmod4$ with $d\ge100$ and every integer $M\ge3$ there is an $L$-family with $B=4dM^2$ such that every prime $p$ with
$\lambda(0,p)\ne0$ is a norm-one prime for $(d,M)$. \textup{(}\lref{GlobalTheory.Hk.exists_lFam_ray},
\lref{GlobalTheory.T.exists_lFam_normOne_of_mod_four_eq_three}.\textup{)}
\end{theorem}

In the formal proof (not converted here), $G$ is a group of ray classes of ideals of $\OO_K$, $K=\Q(\sqrt{-d})$, modulo $M$, $a(g,n)$ is the
number of integral ideals of norm $n$ in the class $g$, $\kappa=2\pi/(\sqrt dM^2)$, and $\lambda(g,n)=\sum_I\Lambda_K(I)$ over the same ideals, where
$\Lambda_K(\mathfrak p^k)=\log\Nm\mathfrak p$ for prime ideals $\mathfrak p$ and $k\ge1$, and $\Lambda_K(I)=0$ otherwise; the
count and theta axioms come from lattice-point counts and theta series of shifted binary quadratic forms
(\lref{GlobalTheory.Hk.rayCountTheta}), the positivity axiom from the non-negativity of the coefficients of products of $L$-series of real
characters (\lref{GlobalTheory.Hk.RaySq.conv_real_char}), and $|G|\le4dM^2$ from a bound for ray class numbers
(\lref{GlobalTheory.Hk.card_ray_le}).

\begin{lemma}[From $\lambda(0,\cdot)$ to norm-one primes]\label{pf:sh-misc:lamnorm}
Let $\mathcal F$ be an $L$-family such that every prime $p$ with $\lambda(0,p)\ne0$ is a norm-one prime for $(d,M)$. Then for all reals $X$ and
$Y\ge1$,
\[
\sum_{\lfloor X\rfloor<n\le\lfloor Y\rfloor}\lambda(0,n)\ \le\ 2\log Y\cdot N_{d,M}(X,Y)+2\bigl(\vartheta_\Lambda(Y)-\vartheta(Y)\bigr),
\]
where $\vartheta_\Lambda(Y)=\sum_{n\le Y}\Lambda(n)$ and $\vartheta(Y)=\sum_{p\le Y}\log p$ are Chebyshev's functions.
\formallabel{GlobalTheory.Hk.sum_lam_zero_le_card_normPrimes}
\end{lemma}

\begin{proof}
By (F5), $\lambda(0,n)\le\sum_g\lambda(g,n)\le2\Lambda(n)$. Split the sum into primes and non-primes. A prime $p$ in the range with
$\lambda(0,p)\ne0$ is a norm-one prime, so the prime part is a sum over at most $N_{d,M}(X,Y)$ primes $p\le Y$ of terms $\le2\log p\le2\log Y$. The
non-prime part is at most $2\sum_{n\le Y,\ n\text{ not prime}}\Lambda(n)=2(\vartheta_\Lambda(Y)-\vartheta(Y))$.
\end{proof}

\begin{lemma}[A numerical lemma]\label{pf:sh-misc:sqrtle}
Let $A,K,c>0$ with $c\le1/(2(1+A)+2+K)$, $B\ge4$, $X\ge1$, $\log B\le c\sqrt{\log X}$ and $e^{-c\sqrt{\log X}}\le h$. Then $\log X\ge1$ and
$K\sqrt X\le hX/B^A$.
\formallabel{GlobalTheory.Hk.sqrt_le_short_main}
\end{lemma}

\begin{proof}
Let $s=\sqrt{\log X}$ and $T=2(1+A)+2+K$. Since $1\le\log B\le cs$ and $cT\le1$, $s\ge1/c\ge T$; so $s\ge1$ and $\log X=s^2\ge1$. As $c\le1$,
$B^A\le e^{As}$ and $h\ge e^{-s}$. Moreover $s^2/2-(1+A)s=s(s/2-(1+A))\ge s\ge K$, so $K\le e^{s^2/2-(1+A)s}$. Hence
$K\sqrt XB^A\le e^{s^2/2-(1+A)s}e^{s^2/2}e^{As}=e^{-s}e^{s^2}\le hX$.
\end{proof}

\begin{proof}[Proof of \cref{pf:sh-misc:shortnorm}]
Let $A,c_0,e_0$ be the constants $A,c,c'$ of \cref{pf:sh-misc:thmA}, and let $C_0$ be such that $\vartheta_\Lambda(x)-\vartheta(x)\le C_0\sqrt x$ for all
$x\ge0$ (Chebyshev; $\vartheta_\Lambda(x)-\vartheta(x)=\sum_{k\ge2}\vartheta(x^{1/k})$). Put $K=8(|C_0|+1)/e_0$,
$c=\min\bigl(c_0,1/(2(1+A)+2+K)\bigr)$ and $c'=e_0/8$. Let $d,M,X,h$ be as in the statement, and let $\mathcal F$ be the $L$-family of
\cref{pf:sh-misc:lfamray}, so $B=4dM^2$. Since $c\le c_0$, \cref{pf:sh-misc:thmA} gives $S:=\sum_{\lfloor X\rfloor<n\le\lfloor X(1+h)\rfloor}\lambda(0,n)\ge
e_0hX/B^A$. By \cref{pf:sh-misc:sqrtle}, $\log X\ge1$ and $8(|C_0|+1)\sqrt X\le e_0hX/B^A$. Let $Y=X(1+h)$; then $1\le Y\le2X$, $\log Y\le\log X+h\le
2\log X$ and $\sqrt Y\le2\sqrt X$. By \cref{pf:sh-misc:lamnorm},
\[
\frac{e_0hX}{B^A}\le S\le4\log X\cdot N_{d,M}(X,Y)+4(|C_0|+1)\sqrt X\le4\log X\cdot N_{d,M}(X,Y)+\frac{e_0hX}{2B^A},
\]
so $N_{d,M}(X,Y)\ge\frac{e_0}8\cdot\frac{hX}{B^A\log X}=\frac{c'hX}{(4dM^2)^A\log X}$.
\end{proof}

